% concatenated from UBmatrix_20260607_.tex
\documentclass[11pt]{article}
\usepackage{amsmath,amssymb,amsthm,verbatim}
\usepackage{color,graphics,srcltx}
\usepackage{hyperref}
\setlength{\oddsidemargin}{-0.15in}
\setlength{\topmargin}{-0.5in}
\setlength{\textheight}{9in}
\setlength{\textwidth}{6.5in}
\newtheorem{theorem}{Theorem}[section]
\newtheorem{definition}[theorem]{Definition}
\newtheorem{proposition}[theorem]{Proposition}
\newtheorem{corollary}[theorem]{Corollary}
\newtheorem{lemma}[theorem]{Lemma}
\newtheorem{example}[theorem]{Example}
\newtheorem{remark}[theorem]{Remark}
\newtheorem{conjecture}[theorem]{Conjecture}
\newtheorem{problem}[theorem]{Problem}
\numberwithin{equation}{section}

\def\rddots{\cdot^{\displaystyle\cdot^{\displaystyle\cdot}}}
\def\IF{{\mathbb F}}
\def\IZ{{\mathbb Z}}
\def\IC{{\mathbb C}}
\def\IR{{\mathbb R}}
\def\IT{{\mathcal T}}
\def\cB{{\mathcal B}}
\def\cC{{\mathcal C}}
\def\cE{{\mathcal E}}
\def\cL{{\mathcal L}}
\def\cR{{\mathcal R}}
\def\cS{{\mathcal S}}
\def\cV{{\mathcal V}}
\def\cP{{\mathcal P}}
\def\cK{{\mathcal K}}
\def\bD{{\mathbb D}}
\def\IH{{\mathbb H}}
\def\bV{{\mathbf V}}
\def\bs{{\mathbf s}}
\def\tr{\mathop{\rm tr}}
\def\span{{\rm span}\,}
\def\la{\langle}
\def\ra{\rangle}
\def\conv{{\rm conv}}
\def\Mat{{\rm Mat}}
\def\bV{{\mathbf V}}
\def\vec{{\rm vec}}
\def\diag{{\rm diag}}
\def\re{{\rm Re\,}}
\def\im{{\rm Im\,}}
\def\rank{{\rm rank\,}}
\def\Ker{\mathop{\mathrm{Ker}}}
\def\char{{\rm char}}
\def\cl{{\bf cl}}
\begin{document}
\openup .88\jot
\title{Numerical radius of certain two-by-two block matrices}
\author{Hwa-Long Gau, Jia-Huo Hong, Chi-Kwong Li, Kuo-Zhong Wang}
\date{}
\maketitle

\centerline{\bf Dedicated to Professor Pei Yuan Wu.}


\begin{abstract}
We investigate the numerical range $W(T)$ and numerical radius $w(T)$ of operators of the form
$T = \begin{pmatrix} A & B \\ 0 & 0 \end{pmatrix}$. We show that $W(T)$ is the union of the
numerical ranges of a family of $2\times 2$ matrices, $T_x$, leading to several consequences,
including improved inequalities for $w(T)$. For cases where $A$ is a self-adjoint involution,
we characterize the conditions under which $W(T)$ is an elliptical disk and determine the minimum
numerical radius of $T_U = \begin{pmatrix} U^*AU & B \\ 0 & 0 \end{pmatrix}$ over all unitary operators
$U$. Finally, we study matrices $T \in M_n$ satisfying $\|T^m x\| = \|T^m\| = \|T\|$ for a unit vector
$x$ and all positive integers $m$. This analysis connects these matrices to the aforementioned block form and
provides a counterexample to the conjecture that if $\|T^k\| = \|T\|$ for all $k \ge 1$, then some power
of the matrix has a direct summand that is a scalar multiple of an idempotent.
\end{abstract}

\noindent \textbf{2020 Mathematics Subject Classification.} Primary 15A60; Secondary 47A12, 15A42.

\noindent \textbf{Keywords.} Numerical range, numerical radius, block matrix, matrix norm, idempotent.

\section{Introduction}




Let $B(H)$ be the algebra of bounded linear operators
acting on the Hilbert space $H$ equipped with the inner product $\la \cdot, \cdot \ra$.
If $H$ has dimension $n$, we identify $B(H)$ as $M_n$,
the algebra of all $n \times n$ complex matrices, and
$H = \IC^n$ equipped with the inner product $\la x, y\ra = y^*x$,
where $y^*$ denotes the conjugate transpose of $y \in \IC^n$.
Let $T \in M_n$. Its operator norm is defined by
\[
\|T\| = \sup \{ \|Tx\| : x \in H,\ \|x\| = 1 \},
\]
where $\|x\|^2 = \langle x, x \rangle$. The numerical range and numerical radius of $T$ are defined
by


\medskip
\centerline{
$W(T) = \{ \langle Tx,x \rangle: x \in H, \|x\| = 1\}$ \quad
and  \quad
$w(T) = \sup \{ |\mu|: \mu \in W(T)\}.$}

\medskip\noindent
The numerical range and numerical radius are useful tools for  studying
operators and matrices; e.g., see \cite{HJ,WG}.
In this paper, we study the numerical range, numerical radius, and some related results
for operator of the form
\begin{equation}\label{T-form}
T = \begin{pmatrix} A & B \\ 0 & 0 \end{pmatrix}.
\end{equation}
Operators of the form \eqref{T-form} appear frequently in both pure and applied areas.
For instance, if $T\in B(H)$ is idempotent, i.e., $T^2 = T$, then $T$ has an operator matrix of
the form (\ref{T-form}) with respect to an orthonormal basis using an orthonormal basis
of the range space of $T$ and its orthogonal complement.  More generally, if $\tilde{T}\in B(H)$
is a quadratic operators, i.e., $\tilde{T}^2 + b\tilde{T} + cI = 0$ for some $b, c \in \IC$,
then $\tilde{T} - \gamma I$ is unitarily similar to an operator $T$ of the form \eqref{T-form}.
By the result in \cite{TW}, $W(\tilde T) - \gamma = W(T)$
is an open or closed elliptical disk with foci at $0$ and $(b^2-4c)^{1/2}$, and with minor axis
of length $\|B\|$. Moreover, $W(T)$ is closed if and only if there
exists a unit vector $x$ such that $\|Bx\| = \|B\|$. There has also been interest in studying
the bounds for $w(T)$ for $T$ in the form \eqref{T-form}; e.g., see \cite{S} and it references.
In section 2, we will show that if $T$ has the form
\eqref{T-form}, then $W(T)$ can be written as the convex hull of elliptical disks; we also
obtain upper bound of $w(T)$ improving some results in the literature. In Section 3, we focus
on $T$ in the form \eqref{T-form} with $A, B \in M_n$ such that $A$ is a Hermitian involution,
i.e., $A = A^*$ and $A^2 = I$.

Another related topic is the
study of matrices $T \in M_n$ satisfying $\|T^m x\| = \|T^m\|=\|T\|$ for some unit vector $x$
and all positive integers $m$ that connects to
matrices of the form \eqref{T-form} with $A \in M_k$ unitary.
In this case, we may always assume that $B \in M_k$. In a sense, it is not surprising as
$T^m = \begin{pmatrix} A^m & A^{m-1}B\cr 0 & 0 \cr\end{pmatrix}$ for $m = 2, 3, \dots$.
As we shall see in Section 4, the study leads to the class of matrices
$\tilde T$ satisfying $\tilde T^*\tilde T = (\tilde T^*)^2 \tilde T^2$.
Evidently, if $\tilde T$ is invertible, then $\tilde T$ satisfies $\tilde T^* \tilde T = I$
so that $\tilde T$ is unitary. If $\tilde T$ is singular, then $\tilde T$ is unitarily similar
to a matrix of form \eqref{T-form}.
In Section 4, we will disprove the following conjecture in \cite{GWW}; see Proposition~4.10.

\begin{conjecture}\label{c11}
If $A \in M_n$ satisfies $\|A^k\| = \|A\|$ for all $k \ge 1$, then some power of $A$ has a direct summand that is a scalar multiple of an idempotent with unimodular scalar.
\end{conjecture}

%Some open problems and possible future research directions will be mentioned in Section 5.




\medskip

We conclude this section with some notations and basic facts about $T$ in the form
\eqref{T-form}.

For $A \in M_n$, let $A^t$ and $A^*$ denote the transpose and conjugate transpose of $A$, respectively.
The singular values of $A$ are the nonnegative square roots of the eigenvalues of $A^*A$. It is known
that $\|A\|$ is the largest singular value of $A$. If $\mathcal{F}\subseteq\{1, \dots, n\}$, we use
$A(\mathcal{F})$ to denote the principal submatrix obtained by deleting the $j$th row and $j$th column
of $A$ for all $j\not\in \mathcal{F}$. We write $\diag(a_1,\ldots,a_n)$ for the diagonal matrix with
diagonal entries $a_1,\ldots,a_n$, and denote by $0_n$ and $I_n$ the zero and identity matrices in
$M_n$, respectively.  We use $0_{k\times m}$ to denote the $k\times m$
zero matrix. For an operator $T\in B(H)$, $T$ is {\em unitarily irreducible} if it is not
unitarily similar to an operator of the form $B\oplus C$; otherwise, it is {\em unitarily reducible}.
We use $\re T$ and $\im T$ to denote the real part $(T+T^*)/2$ and imaginary part $(T-T^*)/(2i)$ of $T$,
respectively. $T$ is {\em positive semidefinite}, denoted by $T\ge 0$, if $\langle Tx,x\rangle\ge 0$ for all $x\in H$. If
$T_1$, $T_2\in B(H)$ are self-adjoint, then $T_1\ge T_2$ means that $T_1-T_2\ge 0$. For any subset $\bigtriangleup$ of complex plane $\IC$, $\cl(\bigtriangleup)$ denotes the closure of $\bigtriangleup$.

In our discussion, it is convenient to embed $H$ into $H\oplus H$ and assume that
$A, B \in B(H)$. We may further replace $T$ by $U^*TU = \begin{pmatrix} U_1^*AU_1 & U_1^*BU_2\cr
0 & 0 \cr\end{pmatrix}$ such that $U_1^*BU_2 \in B(H)$ is positive semidefinite.
In the finite dimensional case, we may assume that $U_1^*BU_2 = \diag(b_1, \dots, b_n)$
with $b_1 \ge \cdots \ge b_n \ge 0$.

\section{Generating ellipses and consequences}

It is known that for any operator $T \in B(H)$, $W(T)$ can be written as the union
of $W(X^*TX)$ for $X: \IC^2 \rightarrow H$ with $X^*X = I_2$.
We will show that if $T$ has the form \eqref{T-form},
then $T$ is the union of $W(X^*TX)$ some special $X$ so that $X^*TX$ has a simple form.
We begin with the following easy case.


\begin{proposition}  Let $T$ have the form \eqref{T-form}. Suppose $A, B \in M_k$ and there are unitary $U, V$ such that
$U^*AU = \diag(a_1, \dots, a_k)$ and $U^*BV = \diag(b_1, \dots, b_k)$.
Let $T_j = \begin{pmatrix} a_j & b_j \cr 0 & 0 \cr\end{pmatrix}$ for $j = 1, \dots, k$.
Then
$$W(T) = \conv \Big(\bigcup_{j=1}^k W(T_j)\Big) \quad \mbox{and} \quad
w(T) = \max\{ w(T_j): 1 \le  j \le k\}.$$
\end{proposition}

\it Proof. \rm This follows from the fact that the assumption ensures that $T$
is a unitarily similar to a direct sum of $T_j$ so that $W(T)=\conv \left(\cup_{j=1}^k W(T_j)\right)$
and $w(T) = \max \{ w(T_j): 1 \le  j \le k \}.$ \qed



\smallskip
In general, we can write $W(T)$ as the union of $W(T_x)$ for some upper triangular
matrix $T_x \in M_2$ even if $T$ does not have the nice special form in the last proposition.

\begin{theorem}\label{p22}
Let $T = \begin{pmatrix} A & B \cr 0 & 0 \cr\end{pmatrix} \in B(H \oplus H)$.
Then
$$W(T) =\conv  \bigcup \{W(T_x): x \in H, \|x\| = 1\} \quad \hbox{ where } T_x
=  \begin{pmatrix}\la Ax,x\ra & \|B^*x\| \cr 0 & 0 \cr\end{pmatrix}.$$
Consequently,
\begin{eqnarray*}
w(T) &=& \sup \{ w(T_x): x \in B(H), \|x\| = 1\} \\
&=& \sup \{ \frac{1}{2}(|\la Ax, x\ra| + \sqrt{|\la Ax,x\ra|^2 + \|B^*x\|^2}): x \in H, \|x\| = 1\}  \\
&\le& \frac{1}{2}w(A) + \frac{1}{2} \min\{ \sqrt{w(AA^*+ BB^*)}, \sqrt{ w(A)^2 + \|B\|^2}\}.
\end{eqnarray*}
The equality holds if and only if there is a sequence of unit vectors $\{x_m\}$
such that $|\langle Ax_m, x_m\rangle| \rightarrow w(A)$ and

{\rm (1)}  $\|B^*x_m\|^2 \rightarrow  w(AA^*+BB^*) - w(A)^2$, or

{\rm (2)} $\|B^*x_m\| \rightarrow  \|B\|$ and $ w(AA^*+BB^*) \ge w(A)^2 + \|B\|^2$.
\end{theorem}

\it Proof. \rm Every unit vector $z \in H\oplus H$ can be written as $\cos \theta x \oplus \sin \theta y$
for some unit vectors $x, y \in H$. Thus, $\la Tz,z\ra = (\cos\theta \ \sin \theta)
\tilde T_x (\cos\theta \ \sin \theta)^t$, where
$\tilde T_x = \begin{pmatrix} \la Ax,x\ra & \la By,  x\ra\cr 0 & 0 \cr\end{pmatrix}$.
Note that $W(\tilde T_x)$ is an elliptical disk with foci $\la Ax,x\ra$ and $0$
with minor axis having length $|\la By,x\ra|$.
We may choose $y = B^*x/\|B^*x\|$ to get $T_x$ as described in the theorem
so that $W(\tilde T_x) \subseteq W(T_x)$. Thus, the set
equality holds.


Now, for any  unit vector $x \in H$,
the numerical radius of $w(T_x)$ is attained at the end point of the major axis and
\begin{eqnarray}\label{eq21}
w(T_x)
 &=&\frac{1}{2}\Big(|\la Ax,x\ra| + \sqrt{|\la Ax,x\ra|^2 + \|B^*x\|^2}\Big) \nonumber\\
&\le&  \frac{1}{2}\Big(|\la Ax,x\ra| + \sqrt{\|A^*x\|^2 + \la BB^*x,x\ra}\Big) \\
&=&  \frac{1}{2}\Big(|\la Ax,x\ra| + \sqrt{\la AA^*x, x\ra + \la BB^*x, x\ra}\Big)  \nonumber\\
&\le&   \frac{1}{2}\Big(w(A) + \sqrt{w(AA^*+BB^*)}\Big).\nonumber
\end{eqnarray}
Clearly, we also have
\begin{equation}\label{eq22}
w(T_x) =\frac{1}{2}\Big(|\la Ax,x\ra| + \sqrt{|\la Ax,x\ra|^2 + \|B^*x\|^2}\Big)
\le \frac{1}{2} \Big(w(A) + \sqrt{w(A)^2+\|B\|^2}\Big).\end{equation}
The equality case could be verified readily by \eqref{eq21} and \eqref{eq22}.
\qed

\medskip
Note that the two quantities $\sqrt{w(AA^*+BB^*)}$ and $\sqrt{w(A)^2 + \|B\|^2}$ are not comparable
in general.
For example,
\begin{itemize}
\item[(a)] if $A = \left({\textstyle 0\atop  \textstyle 0} \ {\textstyle 2 \atop \textstyle 0}\right)$ and $B = \left({\textstyle 1\atop  \textstyle 0} \ {\textstyle 0 \atop \textstyle 0}\right)$, then
$w(AA^* + BB^*) = 5 \ge 2 = w(A)^2 + \|B\|^2$;
\item[(b)] if $A = \left({\textstyle 1\atop  \textstyle 0} \ {\textstyle 0 \atop \textstyle 0}\right)$ and $B=\left({\textstyle 0\atop  \textstyle 0} \ {\textstyle 0 \atop \textstyle 1}\right)$,
then $w(AA^*+BB^*) = 1 \le 2 = w(A)^2 + \|B\|^2$.
 \end{itemize}

\medskip
We have the following corollary, which is  Theorem 3.1 in \cite{S}.
\begin{corollary}  Suppose $T = \begin{pmatrix} A & B \cr 0 & 0 \cr\end{pmatrix} \in B(H \oplus H)$.
Then $w(T) \le \frac{1}{2} (\|A\| + \sqrt{\|AA^* + BB^*\|})$.
\end{corollary}

We also have the following.

\begin{corollary} \label{2.4} Let
$T = \begin{pmatrix} A & B \cr 0 & 0 \cr\end{pmatrix} \in B(H \oplus H)$.
If $B_1, B_2\in B(H)$ such that $B_1B_1^* \le BB^* \le B_2B_2^*$ and $T_j$ is obtained
from $T$ be replacing $B$ with $B_j$ for $j = 1,2$, then
$$W(T_1) \subseteq W(T) \subseteq W(T_2).$$
\end{corollary}

\it Proof. \rm
Note that for any unit vector $x$, $W(T_x)$ is an elliptical disk with foci
$0, \la Ax,x\ra$ and minor axis of length $\|B^*x\|$. For $T_j$ the corresponding
$W((T_j)_x)$ is an elliptical disk with the same foci as $W(T_x)$ and minor axis
of length $\|B_j^*x\|$. Since $B_1B_1^* \le BB^* \le B_2 B_2^*$, we see that
$\|B_1^*x\| \le \|B^*x\| \le \|B_2^*x\|$. Hence,
$W((T_1)_x) \subseteq W(T_x) \subseteq W((T_2)_x)$.
Taking the union of the sets for all unit vectors $x$, we get the inclusion. \qed

\begin{corollary} Let $T$ have the form \eqref{T-form} be such that $B = bI$ with $b \ge 0$ and $W(A)$ is closed.
Then
$$W(T)  = \bigcup \{W(T_x): x\in H,\, \|x\|=1,\,\langle Ax,x\rangle \hbox{ is a boundary point of } W(A)\}.$$
\end{corollary}

\it Proof. \rm Clearly, $W(T_x)$ has the same minor axis for every unit vector $x$.
Hence, the right supporting line of $W(T)$ is determined by the endpoint of the major axis of $W(T_x)$.
Therefore, we only need to focus on those $T_x$ for which $\langle Ax,x\rangle$ is a boundary point of $W(A)$.
\qed

If $A\in M_2$ is unitary, after a rotation,
then $A$ is unitarily similar to $\cos\theta I_2 + i\sin\theta(E_{11}-E_{22})$, where $E_{11}=\left({\textstyle 1\atop  \textstyle 0} \ {\textstyle 0 \atop \textstyle 0}\right)$ and $E_{22}=\left({\textstyle 0\atop  \textstyle 0} \ {\textstyle 0 \atop \textstyle 1}\right)$.
We have the following.

\begin{corollary}
Let $T$ have the form \eqref{T-form} with $A, B \in M_2$ such that
$A = \cos\theta I_2 + i\sin \theta (E_{11}-E_{22})$ with $\theta \in [0,\pi/2]$.
Then $W(T) = \bigcup \{W(T_r): r \in [-\sin\theta , \sin\theta]\}$, where $T_r = \begin{pmatrix} \cos\theta + ir & c \cr 0 & 0 \cr
\end{pmatrix}$, and $(r, c^2)$ ranges over the joint numerical range
$$
W(\im A,\, BB^*) = \{ (\langle (\im A)x,x\rangle,\, \langle BB^* x,x\rangle ) : \|x\|=1 \}.
$$
\end{corollary}



Next, we consider maximum $\sigma(\re (e^{i\theta}T))$, which allows us
to determine  the support lines of $\cl(W(T))$, the closure of $W(T)$, and reduces
the computation of $w(T)$ to an optimization problems over a parameter $\theta \in [0, 2\pi)$.

\begin{theorem} \label{support-line}
Suppose $T\in B(H \oplus H)$ has the form {\rm (\ref{T-form})} with $A, B \in B(H)$.
Let
$$T_\theta = \re(e^{i\theta}T) = \frac{1}{2}
\begin{pmatrix} e^{i\theta} A + e^{-i\theta} A & e^{i\theta} B \cr
e^{-i\theta} B^* & 0 \cr\end{pmatrix}, \qquad \theta \in [0, 2\pi).$$
Then
\begin{eqnarray}\label{eq23}
 \lambda_{\max}(T_\theta) &=&
\inf \{ \lambda \ge 0: 4\lambda\big(\lambda I - \re (e^{i\theta} A)\big)
\ge BB^*\} \nonumber\\
& \ge &   \frac{1}{2} \lambda_{\max}\Big( \re(e^{i\theta}A) +
\sqrt{(\re(e^{i\theta}A))^2 + BB^*}\Big).
\end{eqnarray}
and
\begin{eqnarray*}w(T)&=&\max_{0\le \theta<2\pi}\inf \{ \lambda \ge 0: 4\lambda\big(\lambda I - \re (e^{i\theta} A)\big)
\ge BB^*\}\\
&\ge & \frac{1}{2} \max_{0\le \theta< 2\pi} \lambda_{\max}\Big( \re(e^{i\theta}A) +
\sqrt{(\re(e^{i\theta}A))^2 + BB^*}\Big).\end{eqnarray*}
\end{theorem}

\it Proof. \rm Note that $T_{\theta}$ is
unitarily similar to
$\frac{1}{2}
\begin{pmatrix} e^{i\theta} A + e^{-i\theta} A & B \cr
B^* & 0 \cr\end{pmatrix}$ and $0\in W(T_{\theta})$. Therefore,
$$\lambda_{\max}(T_\theta) = \inf\{\lambda\ge 0 : \lambda I - T_\theta \ge 0\}$$
and
$$\lambda I - T_\theta =  \begin{pmatrix} \lambda I - \re(e^{i\theta} A) & -B/2\cr
-B^*/2 & \lambda I\cr\end{pmatrix},$$
which is positive semidefinite if and only if
$\lambda I - \re(e^{i\theta} A) \ge BB^*/(4\lambda)$. The first equality holds.

Note that $\lambda^2 I - \lambda \re(e^{i\theta} A) -  BB^*/4\ge 0$ is equivalent to
$$\left(\lambda I - \frac{1}{2}\re(e^{i\theta}A)\right)^2
\ge  \frac{1}{4} \left( (\re(e^{i\theta}A))^2 + BB^* \right).$$
Since $\lambda I - T_\theta\ge 0$,
we see that $\lambda I - \re(e^{i\theta}A)\ge 0$ and $\lambda \ge 0$.
So,  $\lambda I - \frac{1}{2}\re(e^{i\theta}A) \ge 0$.
It is  known that if $P \ge  Q \ge 0$ then $\sqrt P \ge \sqrt Q$.
So, $\lambda I - \frac{1}{2}\re(e^{i\theta}A)
\ge \frac{1}{2} \sqrt {(\re(e^{i\theta}  A))^2+BB^*}.$
This shows that
$$\Big\{ \lambda \ge 0:  \lambda I - T_\theta \ge 0\Big\}\subseteq \Big\{ \lambda \ge 0: \lambda I - \frac{1}{2}\Big(\re(e^{i\theta}A)
+ \sqrt {(\re(e^{i\theta}  A))^2+BB^*}\Big)\ge 0\Big\}.$$
That is,
$$\inf \Big\{ \lambda \ge 0:  \lambda I - T_\theta \ge 0\Big\}\ge \inf \Big\{ \lambda \ge 0: \lambda I - \frac{1}{2}\Big(\re(e^{i\theta}A)
+ \sqrt {(\re(e^{i\theta}  A))^2+BB^*}\Big)\ge 0\Big\}.$$
The inequality \eqref{eq23} holds. The final inequality follows from $w(T)=\max_{0\le \theta< 2\pi} \lambda_{\max}(T_{\theta})$.
\qed


\medskip
Note that Theorem \ref{p22} provides upper bounds for $w(T)$ whereas Theorem \ref{support-line} provides a lower bound for $w(T)$. Moreover, in Theorem \ref{p22}, the computation of $w(T)$ reduces to a maximization problem
of $w(T_x)$ for
$x \in H$ with $\|x\| = 1$. In Theorem \ref{support-line}, the computation of $w(T)$ reduced to a
maximization problem
$$\max_{0\le \theta< 2\pi} \inf \{ \lambda \ge 0: 4\lambda\big(\lambda I - \re (e^{i\theta} A)\big)
\ge BB^*\}$$
with two parameters $\lambda \ge 0$ and $\theta \in [0, 2\pi)$. Furthermore,
for a fixed $\theta \in [0, 2\pi)$, one may consider $\lambda > 0$ in a compact set to determine
$$\lambda_\theta = \inf\{\lambda \ge 0:
4\lambda\big(\lambda I - \re (e^{i\theta} A)\big)
\ge BB^*\}$$
and then solve the original optimization by taking $\max\{ \lambda_\theta: \theta \in [0,2\pi]\}$.


\medskip

If $A = \diag(a_1, \dots, a_k)$ and  $B = \diag(b_1, \dots, b_k) \in M_k$, then
\begin{eqnarray*}
\lambda_\theta
&=&  \inf\{ \lambda\ge 0: 4\lambda\big(\lambda - \re(e^{i\theta} a_j) \big)\ge |b_j|^2, j= 1, \dots, k\}\cr
&=&
\frac{1}{2} \max\{ \re(e^{i\theta}a_j) + \sqrt{ (\re(e^{i\theta}a_j))^2 + |b_j|^2}: j = 1, \dots, k\}.
\end{eqnarray*}

A similar formula remains valid for infinite-dimensional operators whenever $\re A$, $\im A$, and $BB^*$ commute, by the spectral theorem.





\section{When the $(1,1)$ block is a self-adjoint involution}

In this section, we consider the special case when
$A$ is a self-adjoint involution operator,
i.e., $A = A^*$ and $A^2= I$.
In such a case, we may assume apply a unitary similar of the form $U = U_1 \oplus U_2$ to $T$,
and assume $A = I_{H_1} \oplus -I_{H_2}$ so that $(\re(e^{i\theta}  A))^2 = \cos^2\theta I$.
By Theorem \ref{support-line}
and the discussion after it,
$$w(T) = \max_{\theta \in [0, 2\pi]}
\inf \{ \lambda \ge 0: 4\lambda\big(\lambda I - \re (e^{i\theta} A)\big)  \ge BB^* \}.$$
In particular, if $A\in M_k$ and $B = \diag(b_1, \dots, b_k)$, then
$$\max\sigma(T_\theta)=
\frac{1}{2}\max\sigma\left(\cos \theta A + \diag\Big(\sqrt{\cos^2\theta +| b_1|^2},
\dots, \sqrt{\cos^2\theta + |b_k|^2}\Big)\right).$$

\iffalse
Also, if $A = A^*$,
then $\sqrt{(\re (e^{i\theta}A))^2 + B^2} = |\re(e^{i\theta}  A)| + B_\theta$ for some
positive semidefinite matrix $B_\theta$.
By Theorem \ref{support-line}, we have
$$\sup \sigma(T_\theta) = \frac{1}{2} \sup  \sigma(\re(e^{i\theta}A) + |\re(e^{i\theta}  A)| + B_\theta).$$
\fi

Suppose  $T$ has the form \eqref{T-form}
such that $A=A^*$ and $A^2 = I$.
We determine the condition on $T$ such that $W(T)$ is an elliptical disk in the following.



\begin{theorem} \label{3.1}
Let  $T=\left(
         \begin{array}{cc}
           A & B \\
           0 & 0 \\
         \end{array}
       \right) \in M_n
$ such that $A = A^* \in M_k$ with $1\le k<n$
satisfying  $A^2 = I_k$ and $B\neq 0$.
The following statements are equivalent.
\begin{itemize}
\item[{\rm (a)}]  $W(T)$ is an elliptical disk.
\item[{\rm (b)}] The matrix $T$ is unitarily similar to a direct sum $T_1\oplus T_2$ with
$W(T_2) \subseteq W(T_1)$, where $T_1$ is either
$$
\left(
  \begin{array}{cc}
    a & \|B\| \\
    0 & 0 \\
  \end{array}
\right)\ \mbox{with } a\in\{1, -1\} \quad \mbox{ or }\quad \left(
                    \begin{array}{ccc}
                      0 & 1 & \|B\| \\
                      1 & 0 & 0 \\
                      0 & 0 & 0 \\
                    \end{array}
                  \right).
$$
\end{itemize}
\end{theorem}

To prove Theorem \ref{3.1}, we need the following lemma.

\begin{lemma}\label{p31}
Let $\theta\in [-\pi,\pi]$, $\alpha_1>\alpha_2\ge 0$,
$$A_\theta=\left(
                                            \begin{array}{cc}
                                              \cos\theta & \sin\theta \\
                                              \sin\theta & -\cos\theta \\
                                            \end{array}
                                          \right),  \ \ \
                                          D=\begin{pmatrix} \alpha_1 & 0 \\ 0 & \alpha_2\end{pmatrix} \ \ \mbox{ and } \ \ T_\theta=\left(
                \begin{array}{cc}
                  A_\theta & D \\
                  0_2 & 0_2 \\
                \end{array}
              \right).$$
Then
$$ w(T_{\pm\pi/2})
=\frac{\sqrt{4+\alpha_1^2+\alpha_2^2+\sqrt{16+8\alpha_1^2+8\alpha_2^2+(\alpha_1^2-\alpha_2^2)^2}}}{2\sqrt{2}} < w(T_\theta) \  \mbox{ for all } \theta\neq \pm \frac{\pi}{2}.
$$
\end{lemma}
\begin{proof}
Observe that $T_\theta$ is unitarily similar to $T_{-\theta}$, so it suffices to consider $\theta\in [0,\pi]$.
For any $\mu\in [0,2\pi]$, $e^{i\mu}T_\theta$ is unitarily similar to
$\left(\begin{array}{cc}
e^{i\mu}A_\theta & D \\
0 & 0 \\
\end{array} \right)$.
Therefore,
\begin{eqnarray*}
&&\mbox{det}(\re(e^{i\mu}T_\theta)-\lambda I_4)=\mbox{det}\left(
                                                                  \begin{array}{cccc}
                                                                    \cos\theta\cos\mu-\lambda & \sin\theta\cos\mu & \frac{\alpha_1}{2} & 0 \\
                                                                    \sin\theta\cos\mu & -\cos\theta\cos\mu-\lambda & 0 & \frac{\alpha_2}{2} \\
                                                                    \frac{\alpha_1}{2} & 0 & -\lambda & 0 \\
                                                                    0 & \frac{\alpha_2}{2} & 0 & -\lambda \\
                                                                  \end{array}
                                                                \right)\\
&=&-\lambda \mbox{det}\left(
                        \begin{array}{ccc}
                          \cos\theta\cos\mu-\lambda& \sin\theta\cos\mu & \frac{\alpha_1}{2} \\
                           \sin\theta\cos\mu &  -\cos\theta\cos\mu-\lambda & 0 \\
                          \frac{\alpha_1}{2} & 0 & -\lambda  \\
                        \end{array}
                      \right)-\frac{\alpha_2^2}{4}\mbox{det}\left(
                                                              \begin{array}{cc}
                                                                \cos\theta\cos\mu-\lambda & \frac{\alpha_1}{2} \\
                                                                \frac{\alpha_1}{2} &  -\lambda \\
                                                              \end{array}
                                                            \right)\\
&=&\lambda^2\mbox{det}\left(
                        \begin{array}{cc}
                         \cos\theta\cos\mu-\lambda& \sin\theta\cos\mu \\
                       \sin\theta\cos\mu &  -\cos\theta\cos\mu-\lambda \\
                        \end{array}
                      \right)-\lambda\frac{\alpha_1^2}{4}(\cos\theta\cos\mu+\lambda)-\frac{\alpha_2^2}{4}(\lambda^2-\lambda \cos\theta\cos\mu-\frac{\alpha_1^2}{4})\\
&=&\lambda^4-\lambda^2(\cos^2\mu+\frac{\alpha_1^2+\alpha_2^2}{4})+\lambda\frac{\alpha_2^2-\alpha_1^2}{4}\cos\theta\cos\mu+\frac{\alpha_1^2\alpha_2^2}{16}\\
&\equiv&F_{\theta,\mu}(\lambda).
\end{eqnarray*}
Let $\lambda_1(F_{\theta,\mu})$ denote the largest real root of $F_{\theta,\mu}$. Then $w(T_\theta)=\max_{\mu\in [0,2\pi]}\lambda_1(F_{\theta,\mu})$.
For $\theta=\pi/2$, since $T_{\pi/2}$ is a nonnegative matrix, we get
\begin{eqnarray*}
\beta & \equiv & w(T_{\pi/2})=\lambda_1(F_{\pi/2,0})=\max\{\lambda :\lambda^4-(1+\frac{\alpha_1^2+\alpha_2^2}{4})\lambda^2+\frac{\alpha_1^2\alpha_2^2}{16}=0\}\\
& = & \frac{\sqrt{4+\alpha_1^2+\alpha_2^2+\sqrt{16+8\alpha_1^2+8\alpha_2^2+(\alpha_1^2-\alpha_2^2)^2}}}{2\sqrt{2}}.
\end{eqnarray*}
For $\theta\in [0,\pi]\setminus \{\pi/2\}$,
$$
F_{\theta,\mu}(\lambda)=F_{\pi/2,0}(\lambda)+\lambda^2\sin^2\mu+\lambda\frac{\alpha_2^2-\alpha_1^2}{4}\cos\theta\cos\mu.$$
Evaluating this at $\lambda=\beta$, we find
$$
F_{\theta,\mu}(\beta)=\beta^2\sin^2\mu-\frac{\alpha_1^2-\alpha_2^2}{4}\beta\cos\theta\cos\mu.
$$
Since $\alpha_1>\alpha_2\ge 0$, we have $F_{\theta, 0}(\beta)<0$ for all $\theta\in[0, \pi/2)$, and, $F_{\theta, \pi}(\beta)<0$ for all $\theta\in(\pi/2, \pi]$, that is, $ \lambda_1(F_{\theta, 0})>\beta$ for all $\theta\in[0, \pi/2)$, and, $ \lambda_1(F_{\theta, \pi})>\beta$ for all $\theta\in(\pi/2, \pi]$. Therefore,
$$w(T_\theta)=\max_{\mu\in [0,2\pi]}\lambda_1(F_{\theta,\mu})\ge \max\{\lambda_1(F_{\theta, 0}), \lambda_1(F_{\theta, \pi})\}>\beta$$
for all  $\theta\in [0,\pi]\setminus \{\pi/2\}$. This completes the proof.
\end{proof}


We are now ready to prove Theorem \ref{3.1}.

\begin{proof}[Proof of Theorem \ref{3.1}]
If $A=\pm I_k$, then $T$ is a quadratic operator and the assertion follows from \cite[Theorem 2.1]{TW}. Therefore, we may assume that $A$ is unitarily similar to $I_{k_1}\oplus (-I_{k_2})$, where $1\le k_1\le k-1$ and $k_2=k-k_1$. Let $m=n-k$, then $B$ is a $k\times m$ matrix. Say, $B=V_0S_B V_1^*$ is the singular value decomposition of $B$, where $V_0\in M_k$ and $V_1\in M_m$ are unitary, and $S_B=[\|B\|]\oplus B_1$ for some $B_1\in M_{k-1, m-1}$. Let $V=V_0\oplus V_1\in M_n$ be unitary, then $V^*TV=\begin{pmatrix} V_0^*AV_0 & S_B\\ 0 & 0\end{pmatrix}$, hence we may assume that $B=[\|B\|]\oplus B_1$ for some $B_1\in M_{k-1, m-1}$.

$(a)\Rightarrow (b)$. Assume that $W(T)$ is an elliptical disk. Since the geometric multiplicity and algebraic multiplicity of each eigenvalue of $T$ are equal, it follows from \cite[Corollary 4]{W} that $W(T)$ is not a circular disk. Therefore, the foci of $W(T)$ must be one of
$\{0,1\}$, $\{0,-1\}$, and $\{-1,1\}$. Moreover, since
$\im T=\left(
                          \begin{array}{cc}
                            0 & B/(2i) \\
                            -B^*/(2i) & 0 \\
                          \end{array}
                        \right)
$, it implies that $\max \sigma (\im T)=\|B\|/2=-\min \sigma (\im T)$, that is, the length of minor axis
of $W(T)$ is $\|B\|$, we infer that either $$W(T)=W(\left(
                                                                 \begin{array}{cc}
                                                                   \pm 1 & \|B\| \\
                                                                   0 & 0 \\
                                                                 \end{array}
                                                               \right)
)\ \ \ \mbox{ or } \ \ \ W(T)= W(\left(
                \begin{array}{cc}
                  1 & \|B\| \\
                  0 & -1 \\
                \end{array}
              \right)
).$$
Since $A^*=A$ and $A^2=I_k$, we may write $A=\left(
                                                     \begin{array}{cc}
                                                       \cos\theta & * \\
                                                       * & A_1 \\
                                                     \end{array}
                                                   \right)
$, where $A_1\in M_{k-1}$ is self-adjoint. By interlacing property, there exists a unitary matrix $U_1\in M_{k-1}$ such that $U_1^*A_1U_1=[\alpha]\oplus I_{k_1-1}\oplus (-I_{k_2-1})$, where $\alpha\in [-1, 1]$. Write $U_1^*B_1=[b_{ij}]\in M_{k-1, m-1}$.
Set $U_2=[1]\oplus U_1\in M_k$ and define $U=U_2\oplus I_m$, so that $U\in M_{n}$ is unitary.  Then $$
U^*TU=\left(
        \begin{array}{cc}
          U_2^*AU_2 & U_2^*B \\
          0 & 0 \\
        \end{array}
      \right)\equiv T',
$$
where
$$U_2^*AU_2=\left(\begin{array}{c|c} \begin{array}{cc} \cos\theta & a_{12}\\ a_{21} & \alpha \end{array} & 0 \\ \hline 0 & \begin{array}{cc} I_{k_1-1} & \\  & -I_{k_2-1}\end{array} \end{array}\right) \ \mbox{ and } \
U_2^*B=\left(\begin{array}{cc|cc} \|B\| & 0 &  \cdots & 0 \\ 0 & b_{11} &  \cdots & b_{1, m-1} \\ \hline \vdots & \vdots &  & \vdots\\ 0 & b_{k-1, 1} & \cdots & b_{k-1, m-1}\end{array}\right).$$
Note that the submatrix $\left({\textstyle \cos\theta \atop  \textstyle a_{21}} \ {\textstyle a_{12} \atop \textstyle \alpha}\right)$ is unitarily similar to $\diag(1, -1)$, thus $\alpha=-\cos\theta$ and $\overline{a}_{21}=a_{12}=e^{i\mu}\sin\theta$ for some real $\mu$, after a unitary similarity by a diagonal unitary matrix, we may assume that $a_{12}=a_{21}=\sin\theta$. On the other hand, the horizontal line $y=\|B\|/2$ is tangent to the elliptical disk $W(T)$ at the point $c+i\|B\|/2$, where $c\in\IR$ is the center of $W(T)$.
Define $x_1=[1\;0\ldots 0]^t\in \IC^k$, $x_2=[i\; 0\ldots 0]^t\in \IC^{m}$, and $x=[x_1\;x_2]^t/\sqrt{2}\in \IC^n$. Then $\|x\|=1$ and $\langle \im(T')x,x\rangle=\|B\|/2$,
so the real part of $\langle T'x,x\rangle$ gives the center of the elliptical disk $W(T)$, or $c=\langle \re(T')x,x\rangle$.
Compute
$$
c=\langle \re(T')x,x\rangle=\langle \left(
                                         \begin{array}{cc}
                                           \cos\theta & \frac{\|B\|}{2} \\
                                           \frac{\|B\|}{2} & 0 \\
                                         \end{array}
                                       \right)\left(
                                                \begin{array}{c}
                                                  \frac{1}{\sqrt{2}} \\
                                                  \frac{i}{\sqrt{2}} \\
                                                \end{array}
                                              \right),
\left(
                                                \begin{array}{c}
                                                  \frac{1}{\sqrt{2}} \\
                                                  \frac{i}{\sqrt{2}} \\
                                                \end{array}
                                              \right)
\rangle=\frac{\cos\theta}{2}.
$$
If $W(T)=W(\left(
             \begin{array}{cc}
               \pm 1 & \|B\| \\
               0 & 0 \\
             \end{array}
           \right)
)$, then $(\cos\theta)/2=c=\pm1/2$ and so, $\cos\theta=\pm 1$.
By interchanging the 2nd and $(k+1)$-th rows and columns of the matrix $T'$, we obtain
$\begin{pmatrix} \pm 1 & \|B\| \cr 0 & 0 \cr \end{pmatrix} \oplus T_2$ as asserted.

Now suppose that $W(T)=W(\left(
                             \begin{array}{cc}
                               1 & \|B\| \\
                               0 & -1 \\
                             \end{array}
                           \right)
)$, then $\cos\theta=2c=0$.
After a unitary similarity by a diagonal unitary matrix, we may assume that the $2\times 2$ leading principal submatrix of $U_2^*AU_2$ is $\left(                                                                                               \begin{array}{cc}                                                                                                   0 & 1 \\                                                                                                   1 & 0 \\                                                                                                 \end{array}                                                                                               \right)$ and $b_{1j}\ge 0$ for all $1\le j\le m-1$.
We now claim that $b_{1j}=0$ for all $1\le j\le m-1$.  Indeed, for each $j=1, \dots, m-1$, consider the $4\times 4$ principal submatrix
$$T'(\{1, 2, k+1, k+1+j\})=
\left(
  \begin{array}{cccc}
    0 & 1 & \|B\| & 0 \\
    1 & 0 & 0 & b_{1j} \\
    0 & 0 & 0 & 0 \\
    0 & 0 & 0 & 0 \\
  \end{array}
\right)\equiv  E_j.$$
We know that $W(E_j)\subseteq W(T')=W(\left(
                        \begin{array}{cc}
                          1 & \|B\| \\
                          0 & -1\\
                        \end{array}
                      \right)
)$, hence
$$w(E_j)\le w(\left(
                        \begin{array}{cc}
                          1 & \|B\| \\
                          0 & -1\\
                        \end{array}
                      \right))=\max \sigma (\left(
                                              \begin{array}{cc}
                                                1 & \frac{\|B\|}{2} \\
                                                \frac{\|B\|}{2} & -1 \\
                                              \end{array}
                                            \right)
                      )=\frac{\sqrt{4+\|B\|^2}}{2}.$$
On the other hand, Lemma \ref{p31} yields that
$$
w^2(E_j)=\frac{4+\|B\|^2+b_{1j}^2+\sqrt{16+8\|B\|^2+8b_{1j}^2+(\|B\|^2-b_{1j}^2)^2}}{8}.
$$
If $b_{1j}\neq 0$, then
$$
w^2(E_j)-\frac{4+\|B\|^2}{4}=\frac{1}{8}\Big(\sqrt{16+8\|B\|^2+8b_{1j}^2+(\|B\|^2-b_{1j}^2)^2}-(4+\|B\|^2-b_{1j}^2)\Big)>0,
$$
a contradiction. We conclude that $b_{1j}=0$ for all $j$.
By interchanging the 3rd and $(k+1)$-th rows and columns of $T'$, we obtain
$$
\left(
  \begin{array}{ccc}
    0 & 1 & \|B\| \\
    1 & 0 & 0 \\
    0 & 0 & 0 \\
  \end{array}
\right)\oplus T_2.
$$

$(b)\Rightarrow (a)$. Let $T_1=\left(
                    \begin{array}{ccc}
                      0 & 1 & \|B\| \\
                      1 & 0 & 0 \\
                      0 & 0 & 0 \\
                    \end{array}
                  \right)$.
We only need to show $W(T_1)$ is an elliptical disk. Since $T_1$ is unitarily similar to
$$
\left(
  \begin{array}{ccc}
    1 & 0 & \|B\|/\sqrt{2} \\
    0 & -1 & \|B\|/\sqrt{2} \\
    0 & 0 & 0 \\
  \end{array}
\right),
$$
by \cite[Theorem 2.2]{KRS}, it is easy to check that the numerical range of the above matrix is an elliptical disk.
This completes the proof.
\end{proof}

It is not hard to extend Theorem \ref{3.1} to compact operators in $B(H)$.
In fact,  one can extend the result to general operators.

\begin{theorem} \label{3.4}
 Suppose $T \in B(H\oplus H)$ has the form \eqref{T-form} such that
$A= A^*$ and $A^2 = I$.
Then  $\cl(W(T))$ is an  elliptical disk ${\mathcal E}$ with foci
in $\{0, -1,1\}$ if and only if  one of
the following holds.
\begin{itemize}
\item[{\rm (a)}] There is a unit vector $x\in H$ such that
$\|B^*x\| = \|B\|$ and $T$ is unitarily similar to $T_1 \oplus T_2$
      with $W(T_2) \subseteq W(T_1)$, where
$$
T_1 = \begin{pmatrix} a & \|B\|\cr 0 & 0\cr\end{pmatrix} \hbox{ with } a \in \{1, -1\} \quad
             \hbox{ or }   \quad
    T_1 = \begin{pmatrix} 0 & 1 & \|B\| \cr  1 & 0 &  0\cr  0 & 0 & 0\cr\end{pmatrix}.$$
\item[{\rm (b)}] There is no unit vectors $x\in H$ such that $\|B^*x\| = \|B\|$,  and there is a sequence of unit
vectors $\{x_m\}$ in $H$ such that one of the following holds.

       \smallskip
        \noindent
        {\rm (b.1)} $\begin{pmatrix} \la Ax_m,x_m\ra &  \|B^*x_m\|\cr  0 & 0 \end{pmatrix}
        \rightarrow  T_0 = \begin{pmatrix} a & \|B\| \cr 0 & 0 \cr \end{pmatrix}$, $a \in \{1, -1\}$,
        and $W(T_0) = \cl (W(T)) = \cE$.

        \smallskip
        \noindent
         {\rm (b.2)}
         $\begin{pmatrix} \la Ax_m,x_m\ra & \sqrt{1-\la Ax_m,x_m\ra^2} &  \|B^*x_m\|\cr   \sqrt{1-\la Ax_m,x_m\ra^2} & -\la Ax_m,x_m\ra & 0 \cr 0 & 0 & 0 \cr \end{pmatrix}
        \rightarrow  T_0 = \begin{pmatrix} 0 & 1 & \|B\| \cr 1 & 0 & 0 \cr 0 & 0 & 0 \cr\end{pmatrix}$
        and

\qquad        $W(T_0) = \cl (W(T)) = \cE$.
\end{itemize}
\end{theorem}

\noindent
\it Proof. \rm
If $A=\pm I$, then $T$ is a quadratic operator and the assertion follows from \cite[Theorem 2.1]{TW}. Therefore, we may assume that $A\neq\pm I$.
The sufficiency proof is clear.
For the necessity, assume that $\cl (W(T))$ is an elliptical disk, where the minor axis
has end points $c+ib/2$, where $b = \|B\|$ and $c \in \{0, \pm 1/2, \pm 1\}$
because the foci are $(0,0), (-1,-1), (1,1), (0,1), (-1,0), (-1,1)$.
Suppose there is a unit vector $x\in H$
such that $\|B^*x\| = \|B\|$.  Consider the $2\times 2$ matrix $T_x=\begin{pmatrix} \langle Ax, x\rangle & \|B^*x\| \\ 0 & 0\end{pmatrix}$. Since $W(T_x)\subseteq W(T)$ and their minor axis have the same length, hence $c=\langle Ax, x\rangle/2$, that is, $|c|\le 1/2$.
It follows that
$c \in\{ 0, \pm 1/2\}$, and hence the foci can only be
$(0,0), (0,1), (-1,0), (-1,1)$.
If the foci are $(0,0)$, then $\cl (W(T))$ is the circular disk centered at 0 with radius
$\|B\|/2$ and thus $\langle Ax, x\rangle=2c=0$.
It implies that $x=(1/\sqrt{2})(u_1+u_2)$, where $u_1\in \ker(A-I)$ and $u_2\in\ker (A+I)$ are unit vectors.
Let $y=(1/\sqrt{2})(u_1-u_2)$,
$\widetilde{x}=x\oplus 0$,
$\widetilde{y}=y\oplus 0$,
and $\widetilde{z}=0\oplus (B^*x/\|B\|)$.
Then $\{\widetilde{x}, \widetilde{y}, \widetilde{z}\}$ is orthonormal and
$$\widetilde{T} \equiv  T|_{\{\widetilde{x}, \widetilde{y}, \widetilde{z}\}} =\begin{pmatrix}
\langle T\widetilde{x}, \widetilde{x}\rangle & \langle T\widetilde{y}, \widetilde{x}\rangle & \langle T\widetilde{z}, \widetilde{x}\rangle\cr
\langle T\widetilde{x}, \widetilde{y}\rangle & \langle T\widetilde{y}, \widetilde{y}\rangle & \langle T\widetilde{z}, \widetilde{y}\rangle\cr
\langle T\widetilde{x}, \widetilde{z}\rangle & \langle T\widetilde{y}, \widetilde{z}\rangle & \langle T\widetilde{z}, \widetilde{z}\rangle\cr
\end{pmatrix}= \begin{pmatrix} 0 & 1 & \|B\|\cr 1 & 0 & 0 \cr 0 & 0 & 0 \cr\end{pmatrix}$$
such that $W(\widetilde{T})$ is not a subset of the circular disk centered at 0 with radius
$\|B\|/2$. Thus, the foci can only be $(0,1), (-1,0), (-1,1)$. By the arguments in the proof of
Theorem \ref{3.1}, we see that $T$ is a unitarily similar to $T_1 \oplus T_2$ as asserted.

Suppose there is no unit vector $x$ such that $\|B^*x\| = \|B\|$.
We can use the Berberian trick in \cite{B} to identify sequences of unit vectors
$\{x_m\}$ in $H$ as an element in
the sequence space. Then identify two sequences $\{x_n\}$ and $\{y_n\}$ if
they differ by a null sequence. We can then use the equivalence classes
in the sequence space to form a Hilbert space
$K$ containing $H$ as the subspace of constant sequences.
Then extend $T$ to an operator
$\tilde T = \begin{pmatrix} \tilde A & \tilde B \cr 0 & 0 \cr\end{pmatrix} \in B(K\oplus K)$,
where $\tilde A, \tilde B\in B(K)$ are the extensions of $A, B\in B(H)$,
i.e., for any sequence $\{x_m\} \in K$, $\tilde A \{x_m\} = \{A x_m\}$ and
$\tilde B \{x_m\} = \{Bx_m\}$. Now, $\tilde B$ will be norm attaining.
By (a),
$\tilde T$ is a direct sum to $\tilde T_0$ and $\tilde T_2$.
Evidently, the norm attaining vector $\{x_m\}$ in $K$ corresponds to
a sequence in $H$ satisfying (b).
\qed

\medskip
\begin{remark}
The necessity of Theorem \ref{3.4} was proved under the assumption that $\cl (W(T))$ is an elliptical disk with foci in $\{0,1,-1\}$.
We conjecture that this assumption is automatic.
\end{remark}

\begin{conjecture}
Let $T\in B(H\oplus H)$ be of the form
\[
T=\begin{pmatrix}A&B\\0&0\end{pmatrix},
\]
where $A=A^*$ and $A^2=I$. If $\cl (W(T))$ is an elliptical disk,
then its foci are contained in $\{0,1,-1\}$.
\end{conjecture}

We are able to prove the conjecture in the case where one of the foci is
$0$; in this case, the other focus must be either $1$ or $-1$.
The general case remains open.






\medskip
 Next we determine the minimum of the numerical radius of $T$ with $A$ such that $A=A^*$ and $A^2=I$
 for a given $B$.  We assume that $A \ne \pm I$ to avoid trivial consideration.
Therefore, we assume that up to unitary similarity, $A = I_p \oplus -I_q$ with
 $p \ge q \ge 1$.  Among other things, suppose $A\in M_k$ and $B\in M_{k, m}$. If $k<m$, by singular value decomposition, we may assume that $B=[B_1 \ 0_{k\times (m-k)}]$ where $B_1\in M_k$, thus $T$ is unitarily similar to $\begin{pmatrix} U^*AU & B_1 \\ 0_k & 0_k\end{pmatrix}\oplus 0_{m-k}$ for some unitary matrix $U\in M_k$. In this case, we have $w(T)=w\Big(\begin{pmatrix} U^*AU & B_1 \\ 0_k & 0_k\end{pmatrix}\Big)$. On the other hand, if $k>m$, let $B'=[B \ 0_{k\times (k-m)}]\in M_k$, then the matrix $T'\equiv\begin{pmatrix} A & B' \\ 0_k & 0_k\end{pmatrix}=T\oplus 0_{k-m}$, and thus $w(T)=w(T')$. Therefore, we only need to consider that both $A$ and $B$ are $k\times k$ matrices.

 For convenience, if $B\in M_k$ has singular values $s_1\ge s_2\ge\cdots\ge s_k\ge 0$, let
\begin{equation}\label{eq1}
S_B=\begin{pmatrix} s_1 & & \\ & \ddots & \\ & & s_k\end{pmatrix}\in M_k \ \ \ \mbox{and} \ \ \ J_k=\begin{pmatrix} & & 1\\ & \rddots & \\ 1 & & \end{pmatrix}\in M_k.
\end{equation}
Among other things, if $A=[a_{ij}]\in M_k$ and $B=[b_{ij}]\in M_k$ satisfy $0\le a_{ij}\le b_{ij}$ for all $1\le i,j\le k$, we write $A\preceq B$. In this case, we have $w(A)\le w(B)$, because a nonnegative matrix attains its numerical radius at a nonnegative unit vector.

Here is our main theorem.

\begin{theorem}\label{t36}
Let $1\le q\le p$, $k=p+q$, $A=I_p\oplus(-I_q)\in M_k$, $B\in M_k$, and
$$\alpha=\min\left\{w\Big(\begin{pmatrix} U^*AU & B\\ 0_k & 0_k\end{pmatrix}\Big) : U\in M_k \mbox{ is unitary}\right\}.$$
Let  $\tilde{ A} = I_p \oplus (-I_p)\in M_{2p}$, $\tilde{B} = B \oplus s_{q+1}I_{p-q}\in M_{2p}$, and
$$\beta=\min\left\{w\Big(\begin{pmatrix} V^* \tilde{A}V & \tilde{B}\\ 0_{2p} & 0_{2p}\end{pmatrix}\Big) : V\in M_{2p} \mbox{ is unitary}\right\}.$$
Then
$$
\alpha=w\Big(\left(\begin{array}{c|c}\begin{array}{ccc}  & & J_q\\ & I_{p-q} & \\ J_q & & \end{array} & S_B
\\ \hline 0_k & 0_k\end{array}\right)\Big)
= w\Big(\begin{pmatrix}
J_{2p} & S_{\tilde{B}} \\ 0_{2p} & 0_{2p}
\end{pmatrix}\Big)=\beta,
$$
where $J_q$, $J_{2p}$, $S_B$, and $S_{\tilde{B}}$ are as in \eqref{eq1}.
\end{theorem}

We need the following lemma to prove Theorem \ref{t36}.

\begin{lemma}\label{l35}
Let $1\le q\le p$, $k=p+q$, $B=\diag(s_1, \dots, s_k)\in M_k$ with $s_1\ge s_2\ge \cdots\ge s_k\ge 0$ and $B_1=\diag(s_{2}, \dots, s_{k-1})\in M_{k-2}$. Then, for any unitary $U\in M_k$, we have
$$
 w\Big(\begin{pmatrix} U^*(I_p\oplus (-I_q))U & B\\ 0_k & 0_k\end{pmatrix}\Big)
\ge   w\Big(\left(\begin{array}{c|c}\begin{array}{cc} 0 & 1\\ 1 & 0\end{array} & \begin{array}{cc} s_1 & 0\\ 0 & s_k\end{array}
\\ \hline 0_2 & 0_2\end{array}\right)\oplus
\begin{pmatrix} {U'}^*(I_{p-1}\oplus (-I_{q-1}))U' & B_1\\ 0_{k-2} & 0_{k-2}\end{pmatrix}\Big) $$
for some unitary matrix $U'\in M_{k-2}$.
\end{lemma}

\begin{proof}
We first show that
$$w\Big(\begin{pmatrix} U^*(I_p\oplus(-I_q))U & B\\ 0_k & 0_k\end{pmatrix}\Big) \ge  w\Big(\left(\begin{array}{c|c}\begin{array}{cc} 0 & 1\\ 1 & 0\end{array} & \begin{array}{cc} s_1 & 0\\ 0 & s_k\end{array}
\\ \hline 0_2 & 0_2\end{array}\right)\Big).$$
Indeed, let $U^*(I_p\oplus(-I_q))U=[a_{ij}]_{i,j=1}^k$ and $A_1=[a_{ij}]_{i,j=2}^k\in M_{k-1}$ be the principal submatrix obtained by deleting the first row and first column of $U^*(I_p\oplus(-I_q))U$. By the interlacing property, there exists a unitary matrix $U_1\in M_{k-1}$ such that $U_1^*A_1U_1=[t_1]\oplus I_{p-1}\oplus (-I_{q-1})$, where $t_1\in [-1, 1]$. Write $U_1^*\diag(s_2, \dots, s_k)U_1=[b_{ij}]_{i,j=1}^{k-1}$, we have $b_{11}\ge s_k$ from the interlacing property. Let $W_1=[1]\oplus U_1\in M_k$ be unitary, then
$$W_1^*(U^*(I_p\oplus(-I_q))U)W_1=\left(\begin{array}{c|c} \begin{array}{cc} a_{11} & a_{12}'\\ a_{21}' & t_1 \end{array} & 0 \\ \hline 0 & \begin{array}{cc} I_{p-1} & \\  & -I_{q-1}\end{array} \end{array}\right)$$
and
$$W_1^*BW_1=\left(\begin{array}{cc|cc} s_1 & 0 &  \cdots & 0 \\ 0 & b_{11} &  \cdots & b_{1, k-1} \\ \hline \vdots & \vdots &  & \vdots\\ 0 & b_{k-1, 1} & \cdots & b_{k-1, k-1}\end{array}\right).$$
Note that the matrix $\left({\textstyle a_{11}\atop  \textstyle a_{21}'} \ {\textstyle a_{12}' \atop \textstyle t_1}\right)$ is unitarily similar to $\diag(1, -1)$ and $b_{11}\ge s_k$, by Lemma \ref{p31} and Corollary \ref{2.4}, we obtain
\begin{eqnarray*}
&& w\Big(\begin{pmatrix} U^*(I_p\oplus(-I_q))U & B\\ 0_k & 0_k\end{pmatrix}\Big)
 \ge    w\Big(\left(\begin{array}{c|c}\begin{array}{cc} a_{11} & a_{12}'  \\ a_{21}' & t_1 \end{array} & \begin{array}{cc} s_1 & 0 \\ 0 & b_{11} \end{array} \\ \hline 0_2 & 0_2\end{array}\right)\Big) \\
&\ge& w\Big(\left(\begin{array}{c|c}\begin{array}{cc} 0 & 1\\ 1 & 0\end{array} & \begin{array}{cc} s_1 & 0\\ 0 & b_{11}\end{array}
\\ \hline 0_2 & 0_2\end{array}\right)\Big)
 \ge  w\Big(\left(\begin{array}{c|c}\begin{array}{cc} 0 & 1\\ 1 & 0\end{array} & \begin{array}{cc} s_1 & 0\\ 0 & s_k\end{array}
\\ \hline 0_2 & 0_2\end{array}\right)\Big)
\end{eqnarray*}
as asserted.

We next show that
$$w\Big(\begin{pmatrix} U^*(I_p\oplus(-I_q))U & B\\ 0_k & 0_k\end{pmatrix}\Big) \ge  w\Big(\begin{pmatrix} {U'}^*(I_{p-1}\oplus (-I_{q-1}))U' & B_1\\ 0_{k-2} & 0_{k-2}\end{pmatrix}\Big)$$
for some unitary matrix $U'\in M_{k-2}$. Indeed, let $A_2=[a_{ij}]_{i,j=1}^{k-1}\in M_{k-1}$ be the leading principal submatrix of $U^*(I_p\oplus(-I_q))U$. By the interlacing property, there exists a unitary matrix $U_2\in M_{k-1}$ such that $U_2^*A_2U_2=[t_2]\oplus I_{p-1}\oplus (-I_{q-1})$, where $t_2\in [-1, 1]$. Write $U_2^*\diag(s_1, \dots, s_{k-1})U_2=[c_{ij}]_{i,j=1}^{k-1}$ and let $C=[c_{ij}]_{i,j=2}^{k-1}\in M_{k-2}$ be the principal submatrix obtained by deleting the first row and first column of $U_2^*\diag(s_1, \dots, s_{k-1})U_2$. Let $W_2=U_2\oplus[1]\in M_k$ be unitary. Then
$$W_2^*(U^*(I_p\oplus(-I_q))U)W_2=\left(\begin{array}{c|c|c}
t_2 & 0 & a_{1k}'' \\ \hline 0 & \begin{array}{cc} I_{p-1} &  \\ & -I_{q-1} \end{array} & 0 \\ \hline a_{k1}'' & 0 & a_{kk}\end{array}\right)$$
and
$$W_2^*BW_2=\left(\begin{array}{c|c|c}
c_{11} & \cdots c_{1, k-1} & 0 \\ \hline
\begin{array}{c} \vdots \\ c_{k-1, 1}\end{array} & C & \begin{array}{c} \vdots \\ 0\end{array} \\ \hline
0 & \cdots \ \ \ 0 & s_k
\end{array}\right),$$
thus $\begin{pmatrix} I_{p-1}\oplus (-I_{q-1}) & C \\ 0_{k-2} & 0_{k-2}\end{pmatrix}$ is a principal submatrix of $\begin{pmatrix} W_2^*(U^*(I_p\oplus(-I_q))U)W_2 & W_2^*BW_2 \\ 0_{k} & 0_{k}\end{pmatrix}$.
Note that $C$ is positive semidefinite with eigenvalues $d_1\ge d_2\ge\cdots\ge d_{k-2}\ge 0$. Let  $D=\diag(d_1, \cdots, d_{k-2})$, then $C=U'D{U'}^*$ for some unitary matrix $U'\in M_{k-2}$. Let $W_3=U'\oplus U'\in M_{2k-4}$ be unitary, we have
$$W_3^*\begin{pmatrix} I_{p-1}\oplus (-I_{q-1}) & C \\ 0_{k-2} & 0_{k-2}\end{pmatrix}W_3=\begin{pmatrix} {U'}^*(I_{p-1}\oplus (-I_{q-1}))U' & D \\ 0_{k-2} & 0_{k-2}\end{pmatrix}.$$ From the interlacing property, we deduce that
$$s_1\ge d_1\ge s_2\ge d_2\ge \cdots\ge s_{k-2}\ge d_{k-2}\ge s_{k-1},$$
thus $D \succeq B_1=\diag(s_2, \dots, s_{k-1})$.  Corollary \ref{2.4} yields that
$$w\Big(\begin{pmatrix} {U'}^*(I_{p-1}\oplus (-I_{q-1}))U' & D \\ 0_{k-2} & 0_{k-2}\end{pmatrix}\Big) \ge w\Big(\begin{pmatrix} {U'}^*(I_{p-1}\oplus (-I_{q-1}))U' & B_1 \\ 0_{k-2} & 0_{k-2}\end{pmatrix}\Big).$$
Therefore, we obtain
\begin{eqnarray*}
& & w\Big(\begin{pmatrix} U^*(I_p\oplus(-I_q))U & B\\ 0_k & 0_k\end{pmatrix}\Big)
\ge w\Big(\begin{pmatrix} I_{p-1}\oplus (-I_{q-1}) & C \\ 0_{k-2} & 0_{k-2}\end{pmatrix}\Big) \\
&=& w\Big(\begin{pmatrix} {U'}^*(I_{p-1}\oplus (-I_{q-1}))U' & D \\ 0_{k-2} & 0_{k-2}\end{pmatrix}\Big)
\ge w\Big(\begin{pmatrix} {U'}^*(I_{p-1}\oplus (-I_{q-1}))U' & B_1\\ 0_{k-2} & 0_{k-2}\end{pmatrix}\Big)\end{eqnarray*}
as asserted. This completes the proof.
\end{proof}

We are now ready to prove Theorem \ref{t36}.

\begin{proof}[Proof of Theorem \ref{t36}]
 First we consider $\alpha$.
Let $B=Q_0S_BQ_1^*$ be the singular value decomposition of $B$ and $W_0=Q_0\oplus Q_1\in M_{2k}$ be unitary, then, for any unitary $U\in M_k$, we have
$$W_0^*\begin{pmatrix} U^*AU & B\\ 0_k & 0_k\end{pmatrix}W_0=\begin{pmatrix} (UQ_0)^*A(UQ_0) & S_B\\ 0_k & 0_k\end{pmatrix}.$$
Since $UQ_0\in M_k$ is also unitary, we deduce that
\begin{equation}\label{eq3}
\alpha=\min\left\{w\Big(\begin{pmatrix} U^*AU & S_B\\ 0_k & 0_k\end{pmatrix}\Big) : U\in M_k \mbox{ is unitary}\right\}\le
w\Big(\left(\begin{array}{c|c}\begin{array}{ccc}  & & J_q\\ & I_{p-q} & \\ J_q & & \end{array} & S_B
\\ \hline 0_k & 0_k\end{array}\right)\Big).
\end{equation}
Therefore, we may assume that $B=S_B=\diag(s_1, \dots, s_k)$. Let $B_j=\diag(s_{j+1}, \dots, s_{k-j})\in M_{k-2j}$ for $j=1, \dots, q$.

Now, for any unitary $U\in M_k$, apply Lemma \ref{l35} repeatedly, we obtain
\begin{eqnarray*}
& & w\Big(\begin{pmatrix} U^*AU & B\\ 0_k & 0_k\end{pmatrix}\Big)   \\
& \ge & w\Big(\left(\begin{array}{c|c}\begin{array}{cc} 0 & 1\\ 1 & 0\end{array} & \begin{array}{cc} s_1 & 0\\ 0 & s_k\end{array}
\\ \hline 0_2 & 0_2\end{array}\right)\oplus
\begin{pmatrix} U_1^*(I_{p-1}\oplus (-I_{q-1}))U_1 & B_1\\ 0_{k-2} & 0_{k-2}\end{pmatrix}\Big) \\
& \ge & w\Big(\left(\begin{array}{c|c}\begin{array}{cc} 0 & 1\\ 1 & 0\end{array} & \begin{array}{cc} s_1 & 0\\ 0 & s_k\end{array}
\\ \hline 0_2 & 0_2\end{array}\right)\oplus
\left(\begin{array}{c|c}\begin{array}{cc} 0 & 1\\ 1 & 0\end{array} & \begin{array}{cc} s_2 & 0\\ 0 & s_{k-1}\end{array}
\\ \hline 0_2 & 0_2\end{array}\right) \oplus
\begin{pmatrix} U_2^*(I_{p-2}\oplus (-I_{q-2}))U_2 & B_2\\ 0_{k-4} & 0_{k-4}\end{pmatrix}\Big) \\
& \ge & \cdots \\
& \ge & w\Big(\left(\begin{array}{c|c}\begin{array}{cc} 0 & 1\\ 1 & 0\end{array} & \begin{array}{cc} s_1 & 0\\ 0 & s_k\end{array}
\\ \hline 0_2 & 0_2\end{array}\right)\oplus \cdots\oplus
\left(\begin{array}{c|c}\begin{array}{cc} 0 & 1\\ 1 & 0\end{array} & \begin{array}{cc} s_q & 0\\ 0 & s_{k-q+1}\end{array}
\\ \hline 0_2 & 0_2\end{array}\right) \oplus
\begin{pmatrix} U_q^*I_{p-q}U_q & B_q\\ 0_{k-2q} & 0_{k-2q}\end{pmatrix}\Big) \\
& = & w\Big(\left(\begin{array}{c|c}\begin{array}{cc} 0 & 1\\ 1 & 0\end{array} & \begin{array}{cc} s_1 & 0\\ 0 & s_k\end{array}
\\ \hline 0_2 & 0_2\end{array}\right)\oplus \cdots\oplus
\left(\begin{array}{c|c}\begin{array}{cc} 0 & 1\\ 1 & 0\end{array} & \begin{array}{cc} s_q & 0\\ 0 & s_{p+1}\end{array}
\\ \hline 0_2 & 0_2\end{array}\right) \oplus
\begin{pmatrix} I_{p-q} & B_q\\ 0_{p-q} & 0_{p-q}\end{pmatrix}\Big),
\end{eqnarray*}
where $U_j\in M_{k-2j}$ is unitary for $j=1, \dots, q$. Note that the matrix
$$\left(\begin{array}{c|c}\begin{array}{ccc}  & & J_q\\ & I_{p-q} & \\ J_q & & \end{array} & S_B
\\ \hline 0_k & 0_k\end{array}\right)$$
is permutationally similar to
$$\left(\begin{array}{c|c}\begin{array}{cc} 0 & 1\\ 1 & 0\end{array} & \begin{array}{cc} s_1 & 0\\ 0 & s_k\end{array}
\\ \hline 0_2 & 0_2\end{array}\right)\oplus \cdots\oplus
\left(\begin{array}{c|c}\begin{array}{cc} 0 & 1\\ 1 & 0\end{array} & \begin{array}{cc} s_q & 0\\ 0 & s_{p+1}\end{array}
\\ \hline 0_2 & 0_2\end{array}\right) \oplus
\begin{pmatrix} I_{p-q} & B_q\\ 0_{p-q} & 0_{p-q}\end{pmatrix},$$
hence we conclude that
$$w\Big(\begin{pmatrix} U^*AU & B\\ 0_k & 0_k\end{pmatrix}\Big) \ge w\Big(\left(\begin{array}{c|c}\begin{array}{ccc}  & & J_q\\ & I_{p-q} & \\ J_q & & \end{array} & S_B
\\ \hline 0_k & 0_k\end{array}\right)\Big)$$
for any unitary $U\in M_k$, that is,
$$\alpha\ge
w\Big(\left(\begin{array}{c|c}\begin{array}{ccc}  & & J_q\\ & I_{p-q} & \\ J_q & & \end{array} & S_B
\\ \hline 0_k & 0_k\end{array}\right)\Big) \ \ \mbox{ or } \ \
\alpha=
w\Big(\left(\begin{array}{c|c}\begin{array}{ccc}  & & J_q\\ & I_{p-q} & \\ J_q & & \end{array} & S_B
\\ \hline 0_k & 0_k\end{array}\right)\Big)$$
by \eqref{eq3}.


Now we turn to $\beta$.
From the proof on $\alpha$, we have
\begin{eqnarray*}
\alpha & = & w\Big(\left(\begin{array}{c|c}\begin{array}{ccc}  & & J_q\\ & I_{p-q} & \\ J_q & & \end{array} & S_B
\\ \hline 0_k & 0_k\end{array}\right)\Big) \\
&= &  w\Big(\left(\begin{array}{c|c}\begin{array}{cc} 0 & 1\\ 1 & 0\end{array} & \begin{array}{cc} s_1 & 0\\ 0 & s_k\end{array}
\\ \hline 0_2 & 0_2\end{array}\right)\oplus \cdots\oplus
\left(\begin{array}{c|c}\begin{array}{cc} 0 & 1\\ 1 & 0\end{array} & \begin{array}{cc} s_q & 0\\ 0 & s_{p+1}\end{array}
\\ \hline 0_2 & 0_2\end{array}\right) \oplus
\begin{pmatrix} I_{p-q} & B_q\\ 0_{p-q} & 0_{p-q}\end{pmatrix}\Big)\\
&=& \max\left\{\max_{1\le j\le q}w\Big(\left(\begin{array}{c|c}\begin{array}{cc} 0 & 1\\ 1 & 0\end{array} & \begin{array}{cc} s_j & 0\\ 0 & s_{k-j+1}\end{array}
\\ \hline 0_2 & 0_2\end{array}\right)\Big), \ \ \max_{q+1\le j\le p} w\Big(\begin{pmatrix} 1 & s_{j}\\ 0 & 0\end{pmatrix}\Big)\right\}.
\end{eqnarray*}
Since $s_{q+1}\ge s_{q+2}\ge \cdots\ge s_p$, we have
$w\Big(\begin{pmatrix} 1 & s_{j}\\ 0 & 0\end{pmatrix}\Big) \le w\Big(\begin{pmatrix} 1 & s_{q+1}\\ 0 & 0\end{pmatrix}\Big)=(1+\sqrt{1+s_{q+1}^2})/2$ for all $q+2\le j\le p$. Therefore, we obtain
\begin{equation}\label{ee1}
\alpha=\max\left\{\max_{1\le j\le q}w\Big(\left(\begin{array}{c|c}\begin{array}{cc} 0 & 1\\ 1 & 0\end{array} & \begin{array}{cc} s_j & 0\\ 0 & s_{k-j+1}\end{array}
\\ \hline 0_2 & 0_2\end{array}\right)\Big), \ \ \frac{1+\sqrt{1+s_{q+1}^2}}{2}\right\}.
\end{equation}
On the other hand, since $\tilde{B} = B \oplus s_{q+1}I_{p-q}$, from \eqref{eq1}, we have
$$S_{\tilde{B}}=\diag(s_1, \dots, s_q)\oplus s_{q+1}I_{p-q}\oplus \diag(\underbrace{s_{q+1}, \dots, s_p}_{p-q})\oplus \diag(\underbrace{s_{p+1}, \dots, s_k}_q).$$
From we have proved before, we obtain
$$
\beta  \equiv  \min\left\{w\Big(\begin{pmatrix} V^* \tilde{A}V & \tilde{B}\\ 0_{2p} & 0_{2p}\end{pmatrix}\Big) : V\in M_{2p} \mbox{ is unitary}\right\}
=  \ w\Big(\left(\begin{array}{cc} J_{2p} & S_{\tilde{B}}
\\ 0_{2p} & 0_{2p}\end{array}\right)\Big). $$
Note that $\left(\begin{array}{cc} J_{2p} & S_{\tilde{B}}
\\ 0_{2p} & 0_{2p}\end{array}\right)$ is permutationally similar to
$$\Bigg(\sum_{j=1}^q\oplus\left(\begin{array}{c|c}\begin{array}{cc} 0 & 1\\ 1 & 0\end{array} & \begin{array}{cc} s_j & 0\\ 0 & s_{k-j+1}\end{array}
\\ \hline 0_2 & 0_2\end{array}\right)\Bigg) \oplus \Bigg(\sum_{j=q+1}^p\oplus\left(\begin{array}{c|c}\begin{array}{cc} 0 & 1\\ 1 & 0\end{array} & \begin{array}{cc} s_{q+1} & 0\\ 0 & s_{j}\end{array}
\\ \hline 0_2 & 0_2\end{array}\right)\Bigg).$$
In particular, the matrix
$$\left(\begin{array}{c|c}\begin{array}{cc} 0 & 1\\ 1 & 0\end{array} & \begin{array}{cc} s_{q+1} & 0\\ 0 & s_{q+1}\end{array}
\\ \hline 0_2 & 0_2\end{array}\right)$$ is unitarily similar to $\begin{pmatrix} 1 & s_{q+1} \\ 0 & 0\end{pmatrix}\oplus \begin{pmatrix} -1 & s_{q+1} \\ 0 & 0\end{pmatrix}$, this yields that
$$w\Big(\left(\begin{array}{c|c}\begin{array}{cc} 0 & 1\\ 1 & 0\end{array} & \begin{array}{cc} s_{q+1} & 0\\ 0 & s_{q+1}\end{array}
\\ \hline 0_2 & 0_2\end{array}\right)\Big)
=w\Big( \begin{pmatrix} 1 & s_{q+1} \\ 0 & 0\end{pmatrix}\Big)
=  \frac{1+\sqrt{1+s_{q+1}^2}}{2}.$$
Moreover, since $s_{q+1}\ge s_{q+2}\ge \cdots\ge s_p$, by Corollary \ref{2.4}, we have
$$w\Big(\left(\begin{array}{c|c}\begin{array}{cc} 0 & 1\\ 1 & 0\end{array} & \begin{array}{cc} s_{q+1} & 0\\ 0 & s_{j}\end{array}
\\ \hline 0_2 & 0_2\end{array}\right)\Big) \le w\Big(\left(\begin{array}{c|c}\begin{array}{cc} 0 & 1\\ 1 & 0\end{array} & \begin{array}{cc} s_{q+1} & 0\\ 0 & s_{q+1}\end{array}
\\ \hline 0_2 & 0_2\end{array}\right)\Big) \ \  \mbox{for all }  q+2\le j\le p.$$
It follows that
$$w\Big(\sum_{j=q+1}^p\oplus\left(\begin{array}{c|c}\begin{array}{cc} 0 & 1\\ 1 & 0\end{array} & \begin{array}{cc} s_{q+1} & 0\\ 0 & s_{j}\end{array}
\\ \hline 0_2 & 0_2\end{array}\right)\Big)= w\Big(\left(\begin{array}{c|c}\begin{array}{cc} 0 & 1\\ 1 & 0\end{array} & \begin{array}{cc} s_{q+1} & 0\\ 0 & s_{q+1}\end{array}
\\ \hline 0_2 & 0_2\end{array}\right)\Big) = \frac{1+\sqrt{1+s_{q+1}^2}}{2}.$$
Hence we obtain
\begin{eqnarray*}
\beta &= & w\Big(\left(\begin{array}{cc} J_{2p} & S_{\tilde{B}}
\\ 0_{2p} & 0_{2p}\end{array}\right)\Big) \\
& = & \max\left\{ w\Big(\sum_{j=1}^q\oplus\left(\begin{array}{c|c}\begin{array}{cc} 0 & 1\\ 1 & 0\end{array} & \begin{array}{cc} s_j & 0\\ 0 & s_{k-j+1}\end{array}
\\ \hline 0_2 & 0_2\end{array}\right)\Big), w\Big( \sum_{j=q+1}^p\oplus\left(\begin{array}{c|c}\begin{array}{cc} 0 & 1\\ 1 & 0\end{array} & \begin{array}{cc} s_{q+1} & 0\\ 0 & s_{j}\end{array}
\\ \hline 0_2 & 0_2\end{array}\right)\Big)  \right\}\\
& = & \max\left\{\max_{1\le j\le q}w\Big(\left(\begin{array}{c|c}\begin{array}{cc} 0 & 1\\ 1 & 0\end{array} & \begin{array}{cc} s_j & 0\\ 0 & s_{k-j+1}\end{array}
\\ \hline 0_2 & 0_2\end{array}\right)\Big), \ \ \frac{1+\sqrt{1+s_{q+1}^2}}{2}\right\} = \alpha
\end{eqnarray*}
from \eqref{ee1}, this completes the proof.
\end{proof}



\medskip
We can extend Theorem \ref{t36} to infinite dimensional operators using the fact
that $B(H)$ is a $C$*-algebra of real rank zero; see \cite{BP}.
The positive semidefinite operator $B$ in its integral representation
$\int_{\lambda} \lambda dE(\lambda)$ can be approximated by a sequence of
 ``diagonal'' operators $\{B_m\}$ each has a finite spectrum.
By a suitable choice of unitary basis, we may consider operators of the form
$T_m = \begin{pmatrix} A_m & B_m \cr 0 & 0 \cr\end{pmatrix}$
with $A_m = \begin{pmatrix} 0 & 0 & J \cr 0 & I & 0 \cr J & 0 & 0 \cr\end{pmatrix}$,
where $J$ is the infinite dimensional extension of $J_k$, say, $J(e_n) = e_{-n}$
if $J \in B(H)$ with $H =\ell_2(\IZ)$, and $J(f(x)) = f(-x)$ if $J \in B(H)$ with $H = L_2(\IR)$.
Then one can show that $w(T_m)$ is minimum among operators of the form
$\begin{pmatrix} UA_mU^* & B_m \cr 0 & 0 \end{pmatrix}$ and  the conclusion also holds
for the limit.

\section{Matrix powers and constant norm attaining vectors}

In this section, we study matrices for which there exists a unit vector
$x$ satisfying
\[
\|A^k x\| = \|A^k\| = \|A\| \  \  \mbox{ for all } k = 1,2,\ldots.
\]
Such a vector will be called a \emph{common norm attaining vector} for the powers of $A$.
This problem is motivated by a simple observation.
If $A \in M_n$ is an idempotent, i.e.,  $A^2 = A$, then the sequence $\{\|A^k\|\}_{k \ge 1}$ is
constant, and there exists
a unit vector attaining these norms simultaneously. This leads naturally to the following question.

\begin{problem} \label{prob}
For which matrices $A \in M_n$ does there exist a unit vector $x$ such that
\[
\|A^k x\| = \|A^k\| = \|A\| \  \  \mbox{ for all } k = 1,2,\ldots?
\]
\end{problem}

Observe that if $\|A\| = \|A^2\| > 0$, then
\[
\|A\| = \|A^2\| \le \|A\|^2 ,
\]
which implies that $\|A\| \ge 1$.

The cases $\|A\| = 1$ and $n = 2$ were studied in \cite{GWW} and \cite{P}, respectively, leading to the following results.
The next two results are from \cite[Corollary 2.6 and Proposition 3.2 (b)]{GWW}.


\begin{proposition}\label{prop41}
For an $n \times n$ matrix $A$, the following are equivalent:
\begin{itemize}
\item[\rm (a)] $\|A^n\| = \|A\| = 1$.
\item[\rm (b)] $\|A^k\| = 1$ for all $k \ge 1$.
\item[\rm (c)] There exists a unit vector $x$ such that $\|A^k x\| = \|A^k\| = \|A\| = 1$ for all $k \ge 1$.
\item[\rm (d)] $A$ is unitarily similar to $[a] \oplus B$ with $\|B\| \le |a| = 1$.
\end{itemize}
\end{proposition}

\begin{proposition}\label{prop42}
For a $2 \times 2$ matrix $A$, the following are equivalent:
\begin{itemize}
\item[\rm (a)] There exists a unit vector $x$ such that $\|A^k x\| = \|A^k\| = \|A\| > 1$ for all $k \ge 1$;
\item[\rm (b)] $A$ is unitarily similar to
\[
\lambda \begin{pmatrix}
1 & b \\
0 & 0
\end{pmatrix},
\]
where $|\lambda| = 1$ and $b \neq 0$.
\end{itemize}
\end{proposition}

Regarding Problem \ref{prob}, the next result shows that the matrix
$T=\begin{pmatrix}
U & B \\
0 & 0
\end{pmatrix}\in M_n$,
where $U$ is unitary, satisfies the condition in Problem \ref{prob}.

\begin{theorem}\label{t410}
Let $T\in M_n$. Then the following are equivalent.
\begin{itemize}
  \item[\rm (a)] $T$ is unitarily similar to a matrix of the form
  \[
  \begin{pmatrix}
    U & B\\
    0 & 0_{n-r}
  \end{pmatrix},
  \]
  where $U\in M_r$, $0\le r\le n$, is unitary.
  \item[\rm (b)] $T^*T={T^k}^*T^k$ for all $k\ge 1$.
  \item[\rm (c)] $T^* T={T^2}^* T^2$.
\end{itemize}
In this case, the following statements hold:
\begin{itemize}
  \item[\rm (i)] $\|T^kx\|=\|T^k\|=\|T\|$ for some unit vector $x\in\IC^n$ and for all $k\ge 1$.
    \item[\rm (ii)] If $T$ is nonzero, then $\|T\|=\sqrt{1+\|B\|^2}$.
  \item[\rm (iii)] If $\lambda$ is a nonzero eigenvalue of $T$, then $|\lambda|=1$.
\end{itemize}
\end{theorem}

\begin{proof}
(b) $\Rightarrow$ (c) is trivial.

(c) $\Rightarrow$ (b). Assume that $T^*T={T^2}^*T^2$. Then, for $k\ge 2$,
$${T^k}^*T^k= {T^{k-2}}^*{T^2}^*T^2T^{k-2}={T^{k-2}}^*T^*TT^{k-2}={T^{k-1}}^*T^{k-1}.$$
Continue this process, we infer that ${T^k}^* T^k={T^{k-1}}^*T^{k-1}={T^{k-2}}^*T^{k-2}= \cdots =T^*T$ as required.

(b) $\Rightarrow$ (a). By the Schur’s triangularization theorem (cf. \cite[Theorem 2.3.1]{HJ}), $T$ is unitarily similar to an upper triangular matrix of the form
$
\begin{pmatrix}
U & B\\
0 & N
\end{pmatrix},
$
where $U$ is an $r \times r$ ($0\le r\le n$) nonsingular upper triangular matrix and $N\in M_{n-r}$ is upper triangular with $\sigma(N)=\{0\}$. We first check that $N=0$. Indeed, since $N$ is nilpotent, then $N^n=0$, and thus
$T^n=
\begin{pmatrix}
U^n & *\\
0 & 0
\end{pmatrix}$. Since $U$ is nonsingular, it follows that
$$
\rank(T^n)=\rank(U^n)=\rank(U).
$$
On the other hand, $\rank(X)=\rank (X^*X)$ for any matrix $X\in M_n$. Therefore, condition (b) yields that
$$\rank(T)=\rank(T^*T)=\rank({T^n}^* T^n)=\rank(T^n)=\rank (U).$$
Hence $\rank(N)=\rank(T)-\rank(U)=0$, or $N=0$ as asserted.

It remains to show that $U$ is unitary for $r\ge 1$. If $r=0$, then $T=N=0$ as asserted form. Therefore, we assume that $r\ge 1$. The equality $T^*T={T^2}^* T^2$ yields that
\[
\begin{pmatrix}
U^*U & U^*B\\
B^*U & B^*B
\end{pmatrix}
=
\begin{pmatrix}
{U^2}^*U^2 & {U^2}^*UB\\
B^*U^*U^2 & B^*U^*UB
\end{pmatrix}.
\]
Comparing the $(1,1)$-blocks yields
$U^*U = {U^2}^*U^2$.
Since $U$ is invertible, we multiply by $(U^*)^{-1}$ on the left and $U^{-1}$ on the right to obtain $I_r = U^*U$.
Hence $U$ is unitary.

(a) $\Rightarrow$ (b). After a unitary similarity, we may assume that
$
T=\begin{pmatrix}
U & B\\
0 & 0_{n-r}
\end{pmatrix},
$
where $U\in M_r$ is unitary.  Let $V=U\oplus I_{n-r}\in M_n$ be unitary, then, for $k\ge 1$, we have
$$T^k=\begin{pmatrix}
U^k & U^{k-1}B \\
0 & 0
\end{pmatrix}=\begin{pmatrix}
U^{k-1} & 0 \\
0 & I_{n-r}
\end{pmatrix}\begin{pmatrix}
U & B \\
0 & 0
\end{pmatrix}=V^{k-1}T,$$
and thus
$${T^k}^*T^k=T^*{V^{k-1}}^*V^{k-1}T=T^*T$$
as required.

Moreover, since $V$ is unitary, it follows that $\|T^k\|=\|V^{k-1}T\|=\|T\|$ for all $k\ge1$. Among
other things, since $T$ is a finite matrix, then $\|Tx\|=\|T\|$ for some unit vector $x\in\IC^n$,
thus $\|T^kx\|=\|V^{k-1}Tx\|=\|Tx\|=\|T\|=\|T^k\|$ for all $k\ge 1$, hence (i) holds.

Next, it is easily seen that
$$
TT^*=\begin{pmatrix}
I_r+BB^* & 0\\
0 & 0
\end{pmatrix},$$
hence $\|T\|^2=\|T^*\|^2=\|TT^*\|=\|I_r+BB^*\|=1+\|BB^*\|=1+\|B^*\|^2=1+\|B\|^2$, this proves (ii).

Finally, (iii) follows from $\sigma(T)\subseteq\sigma(U)\cup\{0\}$ and $U$ is unitary. This completes the proof.
\end{proof}


We aim to formulate a necessary and sufficient condition for a nonzero matrix
$A\in M_n$  satisfying
\[
\|A^k x\| = \|A^k\| = \|A\|  \quad \text{for all } k \ge 1 \mbox{ and for some unit vextor } x\in\IC^n.
\]
The following theorem is the main result of this section, it gives a complete solution of Problem \ref{prob}. Note that the case $\|A\|=1$ has been solved by Proposition \ref{prop41}. Therefore, we need only to consider the case $\|A\|>1$.

\begin{theorem}\label{t44}
Let $A\in M_n$ with $\|A\|>1$, and let $m$ be the degree of the minimal polynomial of $A$.
The following conditions are equivalent:
\begin{itemize}
\item[\rm (a)] There exists a unit vector $x\in \mathbb{C}^n$ such that
\[
\|A^k x\| = \|A^k\| = \|A\|  \quad \mbox{ for all } k = 1,2,\ldots,m+1.
\]

\item[\rm (b)] $A$ is unitarily similar to a matrix of the form
\[
\begin{pmatrix}
T & D\\
0 & C
\end{pmatrix},
\]
where $T\in M_r$, $2 \le r \le m$, satisfies
\[
T^*T = {T^2}^*T^2 \ \ \mbox{ and } \ \
\|A^k\| = \|T^k\| \quad \mbox{ for all }  k = 1,2,\ldots,m+1.
\]
In particular, if $r=2$, then $D=0$.
\end{itemize}
In this case, $A$ is singular.
\end{theorem}



For the proof of Theorem \ref{t44}, we need the next lemma.

In the following, we use $\{e_1, \dots, e_n\}$ to denote the standard basis of $\IC^n$. Moreover, for any matrix $X\in M_n$, if $\|Xx\|=\|X\|$ for some unit vector $x\in\IC^n$, then $X^*Xx=\|X\|^2x$, because $\langle(\|X\|^2I_n-X^*X)x,x\rangle=0$ and $\|X\|^2I_n-X^*X$ is positive semidefinite.


\begin{lemma}\label{l46}
For $n\ge 2$, let $T$ be an $n\times n$ matrix of the form
\[
\begin{pmatrix}
t_{11} & t_{12} & \cdots & t_{1, n-1} & t_{1n}\\
t_{21} & t_{22} & \cdots & t_{2, n-1} & t_{2n}\\
0      & t_{32} & \cdots & t_{3, n-1} & t_{3n}\\
\vdots & \vdots & \ddots & \vdots     & \vdots \\
0      & 0      & \cdots & t_{n, n-1} & t_{nn}
\end{pmatrix},
\]
where $t_{i, i-1}\neq 0$ for $i=2,3,\ldots,n$.
If $\| T^{k}e_{1} \|= \| T^{k} \|= \| T \|$ for $k=1,2,\ldots, n+1$, then
\[
T^*T = {T^2}^*T^{2}.
\]
\end{lemma}

\begin{proof}
Note that $Te_j\in\bigvee\{e_1, \dots, e_{j+1}\}$ for $j=1, \dots, n-1$, it follows that
$$T^2e_j\in T(\bigvee\{e_1, \dots, e_{j+1}\})= \bigvee\{Te_1, \dots, Te_{j+1}\}\subseteq \bigvee\{e_1, \dots, e_{j+2}\}.$$
Continue this process, we have $T^ke_j\in \bigvee\{e_1, \dots, e_{j+k}\}$. In particular,
\begin{equation}\label{eq41}
T^ke_1\in \bigvee\{e_1, \dots, e_{k+1}\} \ \  \mbox{ for } 1\le k\le n-1,
\end{equation}
and thus
\begin{equation}\label{eq42}
\langle T^je_1, e_{k+1}\rangle=0 \ \  \mbox{ for } 0\le j<k\le n-1.
\end{equation}
Moreover, for $k=1,2,\ldots, n-1$, since
$$\langle T^ke_1, e_{k+1}\rangle= \langle T(T^{k-1}e_1), e_{k+1}\rangle = \langle \sum_{j=1}^k\langle T^{k-1}e_1, e_j\rangle Te_j, e_{k+1}\rangle= \langle T^{k-1}e_1, e_k\rangle  \langle Te_k, e_{k+1}\rangle,$$
we infer that
\begin{equation}\label{eq45}
\langle T^k e_1,e_{k+1}\rangle=\langle Te_1,e_{2}\rangle \langle T e_2,e_{3}\rangle \cdots \langle T e_k,e_{k+1}\rangle=t_{21}\cdots t_{k+1, k}\neq 0 \ \ \mbox{ for } k=1,2,\ldots, n-1.
\end{equation}
Therefore, by \eqref{eq41}, \eqref{eq42} and \eqref{eq45}, we obtain successively $e_{k+1}\in \bigvee\{e_1,Te_1,\ldots,T^k e_1\}$ and
$$
\bigvee\{e_1,Te_1,\ldots,T^k e_1\}
=
\bigvee\{e_1,e_2,\ldots,e_{k+1}\}\ \ \mbox{ for } k=0,1,\ldots,n-1.
$$
In particular, this implies that $\{e_1,Te_1, \ldots, T^{n-1}e_1\}$ is a basis of $\mathbb{C}^n$. Therefore, it suffices to prove that
\begin{equation}\label{eq43}
T^*T(T^k e_1) = {T^2}^*T^{2}(T^k e_1)\ \ \mbox{ for all }  k=0,1,\ldots,n-1.
\end{equation}

We proceed by induction on $k$.

\medskip
\noindent
\textbf{(i) The case $k=0$.}
Since $\|Te_1\|=\|T\|=\|T^2\|=\|T^2 e_1\|$, hence
\[
T^*T e_1 = \|T\|^2 e_1 = \|T^2\|^2 e_1 = {T^2}^*T^2 e_1.
\]

\medskip
\noindent
\textbf{(ii) The case $1\le k\le n-1$.}
Assume that \eqref{eq43} holds for $0,\ldots,k-1$, that is,
\begin{equation}\label{eq44}
T^*(T^*T-I_n)T(T^je_1)=0 \ \ \mbox{ for all } j=0,1,\ldots,k-1.
\end{equation}
We want to show that $T^*T(T^k e_1) = {T^2}^*T^{2}(T^k e_1)$. Indeed, Let
$$\mu={T^2}^*T^2(T^k e_1)-T^*T(T^k e_1) = T^*(T^*T-I_n)T(T^ke_1).$$ Then, for $0\le j\le k-1$, \eqref{eq44} yields that
$$\langle \mu, T^je_1\rangle= \langle  T^*(T^*T-I_n)T(T^ke_1), T^je_1\rangle= \langle T^ke_1,  T^*(T^*T-I_n)T(T^je_1)\rangle= \langle T^ke_1, 0\rangle=0.$$
Moreover, by assumption, we have $\|T^{k+2}e_1\|=\|T^{k+2}\|=\|T\|=\|T^{k+1}\|=\|T^{k+1}e_1||$, we infer that
$${T^{k+2}}^*T^{k+2}e_1=\|T^{k+2}\|^2e_1=\|T\|^2e_1=\|T^{k+1}\|^2e_1={T^{k+1}}^*T^{k+1}e_1,$$
or, ${T^{k+1}}^*(T^*T-I_n)T^{k+1}e_1=0$. Therefore, for $k\le j\le n-1$, we obtain
$$\langle \mu, T^je_1\rangle= \langle  T^*(T^*T-I_n)T(T^ke_1), T^je_1\rangle= \langle {T^{k+1}}^*(T^*T-I_n)T^{k+1}e_1, T^{j-k}e_1)\rangle= \langle 0, T^{j-k}e_1 \rangle=0.$$
Since $\{e_1, Te_1, \dots, T^{n-1}e_1\}$ is a basis of $\IC^n$ and $\langle \mu, T^je_1\rangle=0$ for all $0\le j\le n-1$,
we deduce that $\mu=0$, and hence
\[
T^*T (T^k e_1) = {T^2}^*T^2 (T^k e_1).
\]
This completes the proof.
\end{proof}

We are now ready to prove Theorem \ref{t44}.

\begin{proof}[Proof of Theorem \ref{t44}]
(b) $\Rightarrow$ (a).
Assume that
$
A = \begin{pmatrix}
T & D\\
0 & C
\end{pmatrix},
$
where $T$ satisfies the conditions in (b). Theorem \ref{t410} (i) yields that
\[
\|T^kx\|=\|T^k\|=\|T\| \ \ \mbox{ for some unit vector } x\in\IC^r \mbox{ and for all } k\ge 1.
\]
Then
\begin{equation}\label{eq48}
A^k = \begin{pmatrix}
T^k & *\\
0 & C^k
\end{pmatrix}
\quad \text{for all } k\ge 1.
\end{equation}
Set $\widetilde{x} =\left({\textstyle x\atop  \textstyle 0}\right) \in \mathbb{C}^n$. Then $A^k \widetilde{x} = \Big({\textstyle T^kx\atop  \textstyle 0}\Big)$, and for $1 \le k \le m+1$,
\[
\|A^k \widetilde{x}\| = \|T^k x\| = \|T^k\|  = \|A^k\| = \|T\|= \|A\|.
\]
Hence (a) holds.

\medskip

(a) $\Rightarrow$ (b).
Let
\[
K=\bigvee \{x, Ax, \ldots, A^{m-1}x\}\ \ \mbox{ and }  \ \ r = \dim K.
\]
Since $m$ is the degree of the minimal polynomial of $A$, thus $K$ is an invariant subspace of $A$ and $r\le m$. We first show that $r \ge 2$.
If $r=1$, then $Ax = \lambda x$ for some $\lambda\in \mathbb{C}$.
Hence
\[
|\lambda| = \|Ax\| = \|A\| = \|A^2 x\| = |\lambda|^2,
\]
which implies $|\lambda| = 1$, contradicting $\|A\| > 1$.
Thus $r \ge 2$. Since $r=\dim K$,
the vectors $\{x, Ax, \ldots, A^{r-1}x\}$ are linearly independent.
Extend them to an orthonormal basis $\{u_1, u_2, \ldots, u_n\}$ of $\mathbb{C}^n$ such that
\[
\bigvee \{x, Ax, \ldots, A^{k-1}x\}
= \bigvee \{u_1, u_2, \ldots, u_k\}, \quad k=1,\ldots,r,
\]
where $u_1 = x$. Let $U=[u_1 \ \dots \ u_n]\in M_n$ be unitary. Since $K=\bigvee \{x, Ax, \ldots, A^{r-1}x\}=\bigvee \{u_1, u_2, \ldots, u_r\}$ is invariant for $A$, $U^*AU$ has the form
\[
\widetilde{A}\equiv U^*AU = \begin{pmatrix}
T & D\\
0 & C
\end{pmatrix},
\]
where $T=[t_{ij}]\in M_r$ is upper Hessenberg with $t_{i,i-1}\ne 0$ for $i=2,\ldots,r$.
Note that $Ue_1=u_1=x$. Let $e_1' = [1\;0\;\cdots\;0]^t \in \mathbb{C}^r$.
By assumption $\|A^kx\| =\|A^k\|$, $1\le k\le m+1$, we have
\[
\|T^k\| \le  \|\widetilde{A}^k\|=\|A^k\| =\|A^k x\|=\|U\widetilde{A}^kU^*(Ue_1)\|=
\|\widetilde{A}^k e_1\| = \|T^k e_1'\| \le \|T^k\|.
\]
Thus the inequalities above are actually equalities throughout. This implies that
\[
\|A^k\| = \|T^k e_1'\| = \|T^k\|
\]
for $k=1,2,\ldots,m+1$. In particular, $\|T\|=\|A\|=\|A^k\|=\|T^k\|$ for all $1\le k\le m+1$.
That is,
\[
\|T^k e_1'\| = \|T^k\| = \|T\| \quad \text{for } k=1,\ldots,r+1.
\]
By Lemma \ref{l46}, we obtain $T^*T = T^{*2}T^2$
as required.

\medskip

In particular, if $r=2$. Since $T^*T={T^2}^*T^2$, by Theorem \ref{t410} and $\|T\|=\|A\|>1$,
$T$ is unitarily similar to
\[
 \begin{pmatrix}
\alpha & b\\
0 & 0
\end{pmatrix},   \mbox{ where }
 |\alpha|=1 \mbox{ and } b\ne 0.
\]
Hence, without loss of generality, we may assume that $T$ has the above form.
Since $\|\widetilde{A}\| = \|T\| = \sqrt{1+|b|^2}$ and $\|\widetilde{A}^*\| \ge \|\widetilde{A}^* e_1\|$,
we deduce that
\[
D = \begin{pmatrix}
0 & \cdots & 0\\
d_1 & \cdots & d_{n-2}
\end{pmatrix},
\]
and thus the first row of $\widetilde{A}$ is $(\alpha \ \ b \ \ 0 \ \ \dots \ \ 0)$. It follows that the first row of $\widetilde{A}^2$ is
$$(\alpha^2 \ \ \alpha b \ \ bd_1 \ \ \dots \ \ bd_{n-2}).$$
Since $\|\widetilde{A}^2\| =\|\widetilde{A}\|=\sqrt{1+|b|^2}$, $|\alpha|=1$ and $b\neq 0$, it forces that $d_j=0$ for all $1\le j\le n-2$.
Thus $D=0$ as asserted.

Finally, if $A$ is nonsingular, then $T$ is also nonsingular. It follows that
$$I_r=(T^*)^{-1}T^*T(T^{-1})= (T^*)^{-1}{T^2}^*T^2(T^{-1})=T^*T,$$
or $T$ is unitary. Thus $\|A\|=\|T\|=1$, contradicting $\|A\|>1$. Hence $A$ is singular, this completes the proof.
\end{proof}

We have the following consequence immediately.

\begin{corollary}\label{cor412}
Let $A$ be an $n$-by-$n$ nonsingular matrix, and let $m$ be the degree of the minimal polynomial of $A$.
The following conditions are equivalent:
\begin{itemize}
\item[\rm (a)] There exists a unit vector $x \in \mathbb{C}^n$ such that $\|A^k x\| = \|A^k\| = \|A\|$ for all $1\le k \le m+1$.
\item[\rm (b)] $A$ is unitarily similar to $[a] \oplus C$ with $\|C\| \le |a| = 1$.

\item[\rm (c)] There exists a unit vector $x \in \mathbb{C}^n$ such that $\|A^k x\| = \|A^k\| = \|A\|$ for all $k \ge 1$.
\end{itemize}
\end{corollary}

\begin{proof}
The implications (b) $\Rightarrow$ (c) $\Rightarrow$ (a) are trivial. Assume (a) holds, $\|A\|=\|A^2\|\le \|A\|^2$ implies that $\|A\|\ge 1$. On the other hand, since $A$ is nonsingular,  Theorem \ref{t44} yields that $\|A\|\le 1$. Therefore, $1=\|A\|=\|A^k\|$ for all $1\le k\le m+1$, and hence (b) follows from \cite[Theorem 2.1]{P}.
This completes the proof.
\end{proof}

For Theorem \ref{t44}, the next proposition shows that the condition
\[
\|A^k x\| = \|A^k\| = \|A\| > 1 \quad \mbox{for } k=1,2,\ldots,m+1,
\]
cannot, in general, be reduced to $k \le m$.

\begin{proposition}\label{prop48}
For $n \ge 2$, let $A \in M_n$ be of the form
\[
A =
\begin{pmatrix}
0 &        &        &        & 1\\
2 & 0      &        &        &  \\
  & 1      & 0      &        &  \\
  &        & \ddots & \ddots &  \\
  &        &        & 1      & 0
\end{pmatrix}.
\]
Then the following statements hold:
\begin{itemize}
\item[\rm (a)] The minimal polynomial of $A$ has degree $n$.

\item[\rm (b)] $\|A^k e_1\| = \|A^k\| = 2$ for $k=1,2,\ldots,n$.

\item[\rm (c)] $A$ is not unitarily similar to a matrix of the form
\[
\begin{pmatrix}
T & D\\
0 & C
\end{pmatrix},
\]
where $T$ satisfies the condition in Theorem \ref{t44} (b).
\end{itemize}
\end{proposition}

\begin{proof}
(a) A direct computation shows that the characteristic polynomial of $A$ is
\[
\det(xI_n-A) = x^n - 2
= \prod_{k=0}^{n-1} \bigl(x - 2^{1/n} e^{i\frac{2k\pi}{n}}\bigr).
\]
Hence $A$ has $n$ distinct eigenvalues
\begin{equation}\label{eq49}
2^{1/n},\; 2^{1/n} e^{i\frac{2\pi}{n}},\; \ldots,\;
2^{1/n} e^{i\frac{2(n-1)\pi}{n}}.
\end{equation}
It follows that the minimal polynomial of $A$ is $x^n - 2$, and thus has degree $n$.

\medskip

(b) A direct computation yields
\[
A^*A = \operatorname{diag}(4,1,\ldots,1), \quad
{A^2}^*A^2 = \operatorname{diag}(4,1,\ldots,1,4),
\]
\[
\ldots, \quad
{A^{(n-1)}}^*A^{n-1} = \operatorname{diag}(4,1,4,\ldots,4), \quad
{A^n}^*A^n = 4I_n.
\]
Thus
 $\|A^k e_1\| = \|A^k\| = 2$ for $k=1,2,\ldots,n$.

\medskip

(c) By \eqref{eq49}, all eigenvalues of $A$ have modulus $2^{1/n}$.
Thus (c) follows from Theorem \ref{t410} (iii).
\end{proof}


Let $T\in M_r$ satisfy $T^*T={T^2}^*T^2$, and let
\[
A=\begin{pmatrix}
T & D\\
0 & C
\end{pmatrix}.
\]
By Theorem \ref{t410} (b), we have
\[
\|T\|=\|T^k\|\le \|A^k\| \ \ \mbox{ for all } k\ge 1.
\]
Note that in Theorem \ref{t44}, if $r=2$ and $\|A^k\|=\|T\|$ for all $k\ge 1$, then $D=0$. A natural question arises: can Theorem \ref{t44} be extended to the case $r\ge 3$? The answer is negative. A counterexample is given below.

\begin{proposition}\label{prop49}
Let
\[
A=\begin{pmatrix}
1&  0&  1&  0& \frac{1}{4}\\
0&  -1&  1&  0& -\frac{1}{4}\\
0&  0&  0&  0& \frac{1}{4}\\
0&  0&  0&  \frac{1}{4}& \frac{1}{4}\\
0&  0&  0&  0& 0
\end{pmatrix},\quad
T=\begin{pmatrix}
1&  0& 1\\
0&  -1& 1\\
0&  0& 0
\end{pmatrix},
\]
and let $x=\frac{1}{\sqrt{6}}[1 \ -1 \ 2 \ 0 \ 0]^t\in\IC^6$. Then the following hold:
\begin{itemize}
  \item[\rm (a)] $T^*T={T^2}^*T^2$.
  \item[\rm (b)] $\|A^k x\|=\|A^k\|=\|A\|=\|T\|=\sqrt{3}$ for all $k\ge 1$.
  \item[\rm (c)] $A$ is unitarily irreducible.
\end{itemize}
\end{proposition}

\begin{proof}
(a) A direct computation shows that
\[
T^*T=
\begin{pmatrix}
1&  0& 1\\
0&  1& -1\\
1&  -1& 2
\end{pmatrix}
= {T^2}^*T^2.
\]

(b) For $k\ge 1$, one computes
\[
A^{2k}=
\begin{pmatrix}
1&  0&  1&  0& \frac{1}{2}\\
0&  1&  -1&  0& \frac{1}{2}\\
0&  0&  0&  0& 0\\
0&  0&  0&  \frac{1}{4^{2k}}& \frac{1}{4^{2k}}\\
0&  0&  0&  0& 0
\end{pmatrix}
\quad \mbox{and} \quad
A^{2k+1}=
\begin{pmatrix}
1&  0&  1&  0& \frac{1}{2}\\
0&  -1&  1&  0& -\frac{1}{2}\\
0&  0&  0&  0& 0\\
0&  0&  0&  \frac{1}{4^{2k+1}}& \frac{1}{4^{2k+1}}\\
0&  0&  0&  0& 0
\end{pmatrix}.
\]
Moreover,
\[
A^*A=
\begin{pmatrix}
1&  0&  1&  0& \frac{1}{4}\\
0&  1&  -1&  0& \frac{1}{4}\\
1&  -1&  2&  0& 0\\
0&  0&  0&  \frac{1}{16}& \frac{1}{16}\\
\frac{1}{4}&  \frac{1}{4}&  0&  \frac{1}{16}& \frac{1}{4}
\end{pmatrix}.
\]
For $k\ge 1$, one verifies that
\[
{A^k}^*A^k x = 3x
\quad \mbox{and} \quad
T^*T x' = 3x',
\]
where $x'=\frac{1}{\sqrt{6}}[1 \ -1 \ 2]^t\in\IC^3$. Hence $3$ is an eigenvalue of both ${A^k}^*A^k$ and $T^*T$.

Furthermore,
\[
\operatorname{tr}(T^*T)=4<6,
\]
and
\[
\operatorname{tr}({A^k}^*A^k)\le \operatorname{tr}({A^2}^*A^2)
=4+\frac{1}{2}+\frac{2}{4^4}<6.
\]
It follows that $3$ is the largest eigenvalue of both ${A^k}^*A^k$ and $T^*T$. Therefore, we obtain
\[
\|T\|^2=\|T^*T\|=3
\]
and
\[
\|A^k x\|^2
=\langle {A^k}^*A^k x,x\rangle
=3\|x\|^2
=3
=\|{A^k}^*A^k\|
=\|A^k\|^2
\]
for all $k\ge 1$.

(c) It suffices to show that the only projections commuting with $A$ are $0_5$ and $I_5$. Let $P$ be a $5\times 5$ Hermitian projection such that $PA=AP$. Write
\[
P=(p_{ij})_{i,j=1}^5.
\]
Comparing the $(i,j)$-entries of $AP$ and $PA$, we obtain successively
\[
p_{15}=p_{25}=p_{35}=p_{45}=0,
\]
\[
p_{14}=p_{24}=p_{34}=0,\quad p_{55}=p_{44},
\]
\[
p_{13}=p_{23}=0,\quad p_{33}=p_{55},
\]
\[
p_{12}=0,\quad p_{22}=p_{55},
\]
\[
p_{11}=p_{55}.
\]
Hence $P=p_{55}I_5$. Since $P=P^2=P^*$, it follows that $P=0_5$ or $P=I_5$. Therefore, $A$ is unitarily irreducible.
\end{proof}


Finally, we construct a $3$-by-$3$ unitarily irreducible matrix $T$ satisfying
$T^*T = {T^2}^*T^2$, which provides a counterexample to Conjecture \ref{c11}.

\begin{proposition}\label{prop413}
Let
\[
T =
\begin{pmatrix}
e^{i\sqrt{2}\pi} & 0 & \frac{1}{\sqrt{2}}\\
0 & e^{-i\sqrt{2}\pi} & \frac{1}{\sqrt{2}}\\
0 & 0 & 0
\end{pmatrix}.
\]
Then the following statements hold:
\begin{itemize}
\item[\rm (a)] There exists a unit vector $x$ such that
\[
\|T^k x\| = \|T^k\| = \sqrt{2} \quad \text{for all } k \ge 1.
\]

\item[\rm (b)] $T^k$ is unitarily irreducible for all $k \ge 1$.

\item[\rm (c)] For all $k \ge 1$, $T^k$ is not a modulus-one multiple of an idempotent.
\end{itemize}
\end{proposition}

\begin{proof}
(a) In Theorem \ref{t410}, take
\[
U = \begin{pmatrix}
e^{i\sqrt{2}\pi} & 0 \\
0 & e^{-i\sqrt{2}\pi}
\end{pmatrix}
\quad \mbox{and} \quad
B = \begin{pmatrix} \frac{1}{\sqrt{2}} \\[4pt] \frac{1}{\sqrt{2}} \end{pmatrix}.
\]
Since $T= \begin{pmatrix}
U & B \\
0 & 0
\end{pmatrix}$ and $U$ is unitary, the desired conclusion follows from Theorem \ref{t410}.

\medskip

(b) For each $k \ge 1$, the eigenvalues of $T^k$ are
$
e^{ik\sqrt{2}\pi},  e^{-ik\sqrt{2}\pi}$, and $0$.
A direct computation shows that the corresponding eigenvectors can be chosen as
$
e_1,e_2$, and
\[
v = \begin{pmatrix} 1 \\ e^{i2\sqrt{2}\pi} \\ -\sqrt{2}e^{i\sqrt{2}\pi} \end{pmatrix}.
\]
In particular, we have
\[
\langle v, e_1 \rangle \neq 0 \quad \text{and} \quad \langle v, e_2 \rangle \neq 0.
\]
Hence, there is no nontrivial reducing subspace spanned by a subset of eigenvectors, and therefore $T^k$ is unitarily irreducible for all $k \ge 1$.

\medskip

(c) This follows immediately from (b) and \cite[Lemma 3.1]{GWW}.
\end{proof}

Here, \cite[Lemma 3.1]{GWW} says that if $X\in M_n$ is an idempotent and $n\ge 3$, then $X$ is unitarily reducible.

\medskip\noindent
{\bf \large Acknowledgment}

The researches of  Gau and  Wang were partially supported by National Science and Technology Council of Taiwan under grant NSTC 114-2115-M-008-005 and NSTC 114-2115-M-A49-005, respectively.
Li is an affiliate member of the Institute for Quantum Computing, University
of Waterloo. His research was supported by the Simons Foundation Grant 851334.


\begin{thebibliography}{WW}

\bibitem{B}
S.K. Berberian,
Approximate proper values,  Proc.  Amer.  Math.  Soc.  13 (1962), 111--114.

\bibitem{BP}
L.G. Brown and G.K.  Pedersen, $C^*$-algebras of real rank zero,
Journal of Functional Analysis 99 (1991), 131--149.

\bibitem{CL} M.-D. Choi and C.-K. Li,
Numerical ranges and dilations, Linear and Multilinear Algebra 47 (2000),  35--48.

\bibitem{GWW} H.-L. Gau, K.-Z. Wang and Pei Yuan Wu, Constant norms and numerical radii of matrix powers, Operators and Matrices 13 (2019), 1035--1054.

\bibitem{GZ} M. Goldberg and G. Zwas,
On matrices having equal spectral radius and spectral norm,
Linear Algebra Appl. 8 (1974), 427--434.

\bibitem{HJ} R.A. Horn and C.R. Johnson, Matrix Analysis, 2nd ed., Cambridge University Press, Cambridge, 2013.

\bibitem{KRS} D.S. Keeler, L. Rodman, and I.M. Spitkovsky, The numerical range of $3\times 3$ matrices, Linear Algebra Appl. 252 (1977), 115--139.

\bibitem{Li}
C.-K. Li,
Some convexity theorems for the generalized numerical ranges, Linear and Multilinear Algebra 40 (1996), 235--240.

\bibitem{P} V. Pt\'{a}k, Lyapunov equations and Gram matrices, Linear Algebra Appl.   49 (1983), 33--55.

\bibitem{S} K. Shebrawi,
Numerical radius inequalities for certain $2\times 2$
operator matrices II, Linear Algebra Appl.  523 (2017), 1--12.

\bibitem{TW} S.-H. Tso and P. Y. Wu, Matricial ranges of quadratic operators, Rocky Mountain J. Math. 29 (1999), 1139--1152.

\bibitem{W} P.Y. Wu, Numerical ranges as circular discs, Appl. Math. Lett. 24 (2011), 2115--2117.

\bibitem{WG} P.Y. Wu and H.-L. Gau,Numerical Ranges of Hilbert Space Operators, Cambridge University Press, Cambridge, 2021.


\end{thebibliography}


\medskip\noindent
(Gau) Department of Mathematics, National Central University, Chungli 32001, Taiwan. hlgau@math.ncu.edu.tw

\medskip\noindent
(Hong) Department of Mathematics, National Tsing Hua University, Hsinchu 300044, Taiwan. a3229868@gmail.com

\medskip\noindent
(Li) Department of Mathematics, College of William \& Mary, Williamsburg, VA 23187, USA. ckli@math.wm.edu


\medskip\noindent
(Wang) Department of Applied Mathematics, National Yang Ming Chiao Tung University, Hsinchu 300093, Taiwan.
kzwang@math.nctu.edu.tw

\end{document}






